\documentclass[10pt,a4paper]{article}

% ----------------- Paquetes -------------------%
\usepackage[T1]{fontenc}
\usepackage[utf8]{inputenc}
\usepackage[a4paper,margin=2.5cm]{geometry}
\usepackage{mathptmx}             
\usepackage{changepage} 
\usepackage{setspace}
\usepackage[spanish]{babel}
\usepackage[labelfont=bf,labelsep=space]{caption}
\usepackage{amssymb}
\usepackage{amsmath}
\usepackage{ebgaramond}
\usepackage{placeins}
\usepackage{etoolbox}
\usepackage{setspace}
\usepackage{parskip} 
\usepackage{tcolorbox}
\usepackage{needspace}
\usepackage{xparse}
\usepackage{ocgx2}
\usepackage{hyperref}
\usepackage{hypcap}
\usepackage{soul}\sethlcolor{yellow}% resaltado amarillo: \hl{texto} (Panel TCEE, ⌃H)



%%%%%%%%%%%%%%%%%%%%%%%%%%%%%%%%%%%%%%%%%%%%%%%%%%%%%%%%%%%%%%%%
% --- Fija primero tus valores globales "buenos" ---
\setstretch{1.2}                 % Interlineado global
\setlength{\parskip}{0.7\baselineskip}
\setlength{\parindent}{0pt}
\newlength{\GlobalParskip}
\setlength{\GlobalParskip}{\parskip}

\newcounter{eqnum}[subsection]
\renewcommand{\theeqnum}{\arabic{eqnum}}
\newcounter{alphaitem}
\newcounter{numitem}

%% ------------------- Recuadros----------------------%%
\newcommand{\puntosclave}[1]{%
  \begin{tcolorbox}[
    colframe=black,
    title={Puntos clave},
    before upper={\setlength{\parindent}{0.5cm}\setlength{\parskip}{0.5\baselineskip}}
  ]
  #1
  \end{tcolorbox}
}

\newcommand{\modificaciones}[1]{
  \begin{tcolorbox}[
    colback=white,
    colframe=red,
    title={Modificaciones a realizar},
    before upper={\setlength{\parindent}{0.5cm}\setlength{\parskip}{0.5\baselineskip}}
  ]
  #1
  \end{tcolorbox}
}

%% ------------------- Preguntas TEST----------------------%%
% Opciones: radio (single-choice)
\NewDocumentCommand{\OptRadio}{m m m}{%
  \noindent\CheckBox[radio,name=#1,value=#2,width=1.2ex,height=1.2ex]{}%
  \;\textbf{#2.}~#3\par
}

% Opciones: checkbox (multi-choice)
\NewDocumentCommand{\OptCheck}{m m}{%
  \noindent\CheckBox[name=#1,width=1.2ex,height=1.2ex,bordercolor={0 0 0}]{}%
  \;#2\par
}

% Pregunta single-choice (radio)
% Args: 1=id, 2=enunciado, 3=opciones, 4=solución+justificación, 5=imagen (o {}), 6=origen
\NewDocumentCommand{\PreguntaTestSingle}{m m m m m m}{%
  \par\medskip
  \noindent\textbf{#1.}~#2\par\smallskip

    % Imagen (si no es {})
    \IfBlankTF{#5}{}{%
      \begin{center}
        \includegraphics[width=0.75\linewidth]{#5}
      \end{center}
    }

    \begin{Form}
      #3
    \end{Form}

    \par\smallskip
    \switchocg{sol-#1}{\fbox{Mostrar / ocultar solución}}%
    \hfill{\footnotesize #6}\par\smallskip

    \begin{ocg}{Solución}{sol-#1}{0}
      \noindent #4
    \end{ocg}
  \par\medskip
}

% Pregunta multi-choice (checkboxes)
\NewDocumentCommand{\PreguntaTestMulti}{m m m m m m}{%
  \par\medskip
  \noindent\textbf{#1.}~#2\par\smallskip

    % Imagen (si no es {})
    \IfBlankTF{#5}{}{%
      \begin{center}
        \includegraphics[width=0.75\linewidth]{#5}
      \end{center}
    }
    \begin{Form}
      #3
    \end{Form}

    \par\smallskip
    \switchocg{sol-#1}{\fbox{Mostrar / ocultar solución}}%
    \hfill{\footnotesize #6}\par\smallskip

    \begin{ocg}{Solución}{sol-#1}{0}
      \noindent #4
    \end{ocg}
  \par\medskip
}

%% ------------------- Listas----------------------%%
    % --- Lista alfabética ---
    \newenvironment{la}
      {\begin{list}{\alph{alphaitem})}{%
          \usecounter{alphaitem}%
          \setlength{\leftmargin}{1cm}%
          \setlength{\parskip}{3pt}% 
          \setlength{\itemsep}{0pt}% ← espacio entre ítems
          \setlength{\parsep}{0pt}%  ← párrafos dentro de un ítem
      }}
      {\end{list}}
      
    % --- Lista numérica ---
    \newenvironment{lnum}
      {\begin{list}{\arabic{numitem}.}{%
          \usecounter{numitem}%
          \setlength{\leftmargin}{1cm}%
          \setlength{\parskip}{3pt}% 
          \setlength{\itemsep}{0pt}% ← espacio entre ítems
          \setlength{\parsep}{0pt}%  ← párrafos dentro de un ítem
      }}
      {\end{list}}

%% ------------------- Preguntas Test ----------------------%%
      \usepackage{xparse} % si ya lo tienes, no lo repitas

    \NewDocumentEnvironment{test}{m m o}{%
      \begin{tcolorbox}[
        colback=white,
        colframe=black!15,
        boxrule=0.6pt,
        arc=2mm,
        left=2mm,right=2mm,top=2mm,bottom=2mm
      ]
      \centering
      \includegraphics[width=\linewidth]{#1}
      \end{tcolorbox}
    
      \vspace{2mm}
    
      % Caja SOLO para texto (SÍ partible y mantiene márgenes al partir)
      \begin{tcolorbox}[
        enhanced,
        breakable,
        colback=black!3,
        colframe=black!25,
        boxrule=0.6pt,
        arc=2mm,
        left=2mm,right=2mm,top=1.5mm,bottom=1.5mm
      ]
      \textbf{Respuesta:} \textbf{#2}%
      \IfValueT{#3}{\par\smallskip \textbf{Justificación:} #3}
      \end{tcolorbox}
    }{}
      
% ------------------- Imágenes----------------------%
    \graphicspath{{Img/}}% Rutas por defecto 
    % --- Imagen adaptativa ---
    % Registros de dimensión definidos una sola vez
    \newdimen\imgcap
    \newdimen\imgmin
    \newdimen\imgmax
    \newdimen\imgremain
    \newdimen\imgh
    \newdimen\imgneed
    
    \newcommand{\imagenfit}[3]{%
      \addvspace{10pt}% separación antes
      \par\begingroup
        % Parámetros de composición (asignaciones, no \newdimen)
        \imgcap=1\baselineskip            % alto estimado de título+fuente
        \imgmin=0.15\textheight
        \imgmax=0.15\textheight
        \imgremain=\pagegoal
        \ifdim\pagegoal=\maxdimen \imgremain=0pt\fi
        \advance\imgremain by -\pagetotal
        \advance\imgremain by -\imgcap
        % Decisión de altura destino
        \ifdim\imgremain<\imgmin
          \newpage
          \imgh=0.20\textheight
        \else
          \imgh=\imgremain
          \ifdim\imgh>\imgmax \imgh=\imgmax \fi
        \fi
        % Reservar espacio para evitar cortes del bloque completo
        \imgneed=\imgh \advance\imgneed by \imgcap
        \Needspace{\imgneed}
        % Cuerpo
        \noindent\begin{minipage}{\textwidth}
          \centering
          \captionof{figure}{\textemdash\ #2}\par\vspace{0.3em}% título arriba
          \includegraphics[width=\linewidth,height=\imgh,keepaspectratio]{\detokenize{#1}}\par
          \vspace{0.3em}\small\textit{Fuente: }#3% fuente abajo
        \end{minipage}%
      \endgroup
      \par\addvspace{10pt}%
    }

%------------ Bloque de RECUADRO ESPECIAL -------------%
    \newenvironment{specialblock}[1]{%
      \begin{quote} % crea sangría a izquierda y derecha
      \small % texto más pequeño
      \noindent\textbf{#1}\\[-0.3em] % título en negrita
      \rule{\linewidth}{0.4pt}\\[0.6em]}{
      \end{quote}}
      
%-------------------------- Portada -----------------------%
\newcommand{\oldtitle}[1]{\def\OldTitle{#1}}
\newcommand{\changerationale}[1]{\def\ChangeWhy{#1}}

\newcommand{\changebox}{%
\begin{tcolorbox}[title={Historial de cambios del tema}]
\textbf{Título anterior:} \OldTitle\par
\textbf{Motivo del cambio:} \ChangeWhy
\end{tcolorbox}
}

%-------------------------- Author Font -----------------------%
    \newcommand{\authorfont}[1]{{\fontfamily{EBGaramond-TLF}\selectfont #1}}

%-------------------------- Anexos -----------------------%
    \makeatletter
    \newcounter{anexo}
    \renewcommand{\theanexo}{\Roman{anexo}}
    \newcommand{\anexo}[1]{%
      \clearpage
      \phantomsection
      \refstepcounter{anexo}%
      \noindent\textbf{Anexo \theanexo.- }#1\par\medskip
      \addcontentsline{stc}{section}{Anexo \theanexo.- #1}%
    }
    \makeatother

    %-------------------------- Lista de figuras (personal) -----------------------%
\makeatletter
\newcommand{\MyListOfFigures}{%
  \clearpage
  \phantomsection
  {\centering\bfseries\Large Índice de figuras\par}%
  \medskip
  % Entrada en nuestro índice de contenido (stc)
  \addcontentsline{stc}{section}{Índice de figuras}%
  % Recuperación del listado (sin cabecera propia)
  \begingroup
    \parindent=0pt\parskip=2pt
    \@starttoc{lof}%
  \endgroup
}
\makeatother

%-------------------------- Bibliografía -----------------------%
\makeatletter
% Contador para las referencias
\newcounter{myref}
% Cabecera de la bibliografía, entra en el índice 'stc'
\newcommand{\MyBibliography}{%
  \clearpage
  \phantomsection
  {\centering\bfseries\Large Bibliografía y referencias\par}%
  \medskip
  \addcontentsline{stc}{section}{Bibliografía y referencias}%
  \setcounter{myref}{0}%
}
\newcommand{\MyBibItem}[4]{%
  \refstepcounter{myref}%
  \noindent[\themyref]\ #1.\ \emph{#2}. #3. #4\par
}
\makeatother

%--------- Establecer cuatro niveles de índice --------%
    \setcounter{tocdepth}{4}   
    \setcounter{secnumdepth}{4}% numera hasta \paragraph

%%%%%%%%%%%%%%%%%%%%%%%%%%%%%%%%

\makeatletter

    %--------- Interlineado Global --------%
    \edef\GlobalBaseLineStretch{\baselinestretch}
    
    %--------- Espaciado de seccinones --------%
    \renewcommand\section{\@startsection{section}
      {1}{\z@}
      {2.5ex \@plus 1ex \@minus .2ex} % espacio antes
      {1.5ex \@plus .2ex}             % espacio después
      {\normalfont\Large\bfseries}    % formato del título
    }
    \renewcommand\subsection{\@startsection{subsection}{2}{\z@}
      {3ex \@plus 0.8ex \@minus .2ex}
      {1.5ex \@plus .2ex}
      {\normalfont\large\bfseries}
    }
    \renewcommand\subsubsection{\@startsection{subsubsection}{3}{\z@}
      {3ex \@plus 0.5ex \@minus .2ex}
      {1ex \@plus .2ex}
      {\normalfont\normalsize\bfseries}
    }
    \renewcommand\paragraph{\@startsection{paragraph}{4}{\z@}%
    {-2ex \@plus -0.5ex \@minus -.2ex}% espacio antes
    {0.8ex \@plus .2ex}% espacio después
    {\normalfont\normalsize\itshape}}% ← cursiva en lugar de negrita

    %--------- Ecuaciones --------%   
    % Variables globales para almacenar títulos y notas
    \gdef\leq@current@section{}%
    \gdef\leq@current@subsection{}%
    \gdef\leq@last@printed@section{}%
    \gdef\leq@last@printed@subsection{}%
    
    % Comando para añadir nota (invisible en el cuerpo)
    \newcommand{\eqnote}[1]{%
      \addcontentsline{leq}{leqnote}{#1}%
    }
    
    % Comando para ecuaciones
    \newcommand{\eqblock}[2]{%
      % Verificar si necesitamos imprimir el título de sección
      \ifx\leq@current@section\leq@last@printed@section\else
        \addcontentsline{leq}{leqsection}{\leq@current@section}%
        \global\let\leq@last@printed@section\leq@current@section
        \gdef\leq@last@printed@subsection{}% Reset subsección al cambiar sección
      \fi
      % Verificar si necesitamos imprimir el título de subsección
      \ifx\leq@current@subsection\leq@last@printed@subsection\else
        \ifx\leq@current@subsection\@empty\else
          \addcontentsline{leq}{leqsubsection}{\leq@current@subsection}%
          \global\let\leq@last@printed@subsection\leq@current@subsection
        \fi
      \fi
      % Ecuación
      \refstepcounter{eqnum}%
      \begin{equation*}
        #1\tag{E.\theeqnum}
      \end{equation*}%
      \addcontentsline{leq}{leqt}{[E.\theeqnum] -- #2}%
      \addcontentsline{leq}{leqm}{#1}%
    }
    
    %%% ----- Índice de ecuaciones ----------%%%
    % Tamaño para []
    \newcommand{\leqIndexFont}{\normalsize}
    \newcommand{\leqTitleFont}{\tiny}
    \newcommand{\leqNoteFont}{\tiny}
    % Cabecera de sección 
    \newcommand*\l@leqsection[2]{%
      \addvspace{25pt}% Mayor separación antes de nueva sección
      \noindent\leqIndexFont
      \makebox[\linewidth]{\rule{.38\linewidth}{0.4pt}\ \ \textbf{  \MakeUppercase{#1}}\ \ \rule{.38\linewidth}{0.4pt}}%
      \par\vspace{-10pt}
    }
    % Cabecera de subsección 
    \newcommand*\l@leqsubsection[2]{%
      \par\addvspace{15pt}%
      \noindent\leqIndexFont #1
      \par\vspace{-7pt}
    }
    % Título de ecuación 
    \newcommand*\l@leqt[2]{%
    \par\addvspace{10pt}%
      \par\noindent\leqTitleFont #1\par
    }
    % Miniatura de ecuación
    \newcommand*\l@leqm[2]{%
      \vspace{0.3\baselineskip}%
      \par\scriptsize\noindent\hspace*{1.5em}\(#1\)\par%
    }
    % Nota de ecuación 
    \newcommand*\l@leqnote[2]{%
      \vspace{0.2\baselineskip}%
      \par\noindent\leqNoteFont\hspace*{1.5em}\textbf{Nota.--} #1\par\vspace{0.5\baselineskip}%
    }
    % Generación del índice
    \newcommand{\mylistofequations}{%
      \clearpage
      \phantomsection
      % Título visual de la lista (ajústalo a tu gusto):
      {\centering\bfseries\Large Índice de ecuaciones\par}
      % Entrada en nuestro índice 'stc':
      \addcontentsline{stc}{section}{Índice de ecuaciones}%
      % Llamada a tu lista real (usa el nombre exacto de tu macro):
      \@starttoc{leq}%
    }
    
    %--------- Captura de títulos de sección --------%
    \let\leq@original@section\section
    \let\leq@original@subsection\subsection
    % Flag para saber si estamos en una sección asterisco
    \newif\ifLeqStarSection
    \LeqStarSectionfalse
    
    % Redefinir \section CON CONTROL DE ASTERISCO
    \renewcommand{\section}{%
      \@ifstar{\leq@section@star}{\@dblarg\leq@section@nostar}%
    }
    \def\leq@section@star#1{%
      \global\LeqStarSectiontrue
      \gdef\leq@current@section{#1}%
      \gdef\leq@current@subsection{}%
      \leq@original@section*{#1}%
      \global\LeqStarSectionfalse
    }
    \def\leq@section@nostar[#1]#2{%
      \global\LeqStarSectionfalse
      \gdef\leq@current@section{#2}%
      \gdef\leq@current@subsection{}%
      \leq@original@section[#1]{#2}%
    }
    % Redefinir \subsection
    \renewcommand{\subsection}{%
      \@ifstar{\leq@subsection@star}{\@dblarg\leq@subsection@nostar}%
    }
    \def\leq@subsection@star#1{%
      \global\LeqStarSectiontrue
      \gdef\leq@current@subsection{#1}%
      \leq@original@subsection*{#1}%
      \global\LeqStarSectionfalse
    }
    \def\leq@subsection@nostar[#1]#2{%
      \global\LeqStarSectionfalse
      \gdef\leq@current@subsection{#2}%
      \leq@original@subsection[#1]{#2}%
    }

    % ----- Restauración de valores Globales de estilo ----------%
    % Flags
    \newif\ifInFootnote
    \InFootnotefalse
    \pretocmd{\@footnotetext}{\InFootnotetrue}{}{}
    \apptocmd{\@footnotetext}{\InFootnotefalse}{}{}
    % Profundidad de listas (soporta anidadas)
    \newcount\MyListDepth \MyListDepth=0
    \AtBeginEnvironment{la}{\advance\MyListDepth by 1}
    \AfterEndEnvironment{la}{\advance\MyListDepth by -1}
    \AtBeginEnvironment{lnum}{\advance\MyListDepth by 1}
    \AfterEndEnvironment{lnum}{\advance\MyListDepth by -1}
    \AtBeginEnvironment{itemize}{\advance\MyListDepth by 1}
    \AfterEndEnvironment{itemize}{\advance\MyListDepth by -1}
    \AtBeginEnvironment{enumerate}{\advance\MyListDepth by 1}
    \AfterEndEnvironment{enumerate}{\advance\MyListDepth by -1}
    % Restauración diferida: hazlo al INICIO del próximo párrafo real
    \newif\if@pendingRestore
    \@pendingRestorefalse
    \newcommand{\RequestRestore}{\global\@pendingRestoretrue}
    \everypar{%
      \if@pendingRestore
        \global\@pendingRestorefalse
        \ifInFootnote\else
          \ifnum\MyListDepth=0
            \setlength{\parskip}{\GlobalParskip}%
            \setstretch{\GlobalBaseLineStretch}%
          \fi
        \fi
      \fi
    }
    % Al final de bloques:
    \AfterEndEnvironment{equation}{\RequestRestore}
    \AfterEndEnvironment{equation*}{\RequestRestore}
    \AfterEndEnvironment{la}{\RequestRestore}
    \AfterEndEnvironment{lnum}{\RequestRestore}
    % Al final de macros:
    \apptocmd{\eqblock}{\RequestRestore}{}{}
    \apptocmd{\imagenfit}{\RequestRestore}{}{}
    
\makeatother

% ===================== ToC propio con 4 niveles (único bloque) =====================
\makeatletter

% --- Escritores: añadimos una línea al .stc en cada cabecera numerada ---
% Guardamos las versiones actuales para llamar al comportamiento normal
\let\MyTOC@orig@section\section
\let\MyTOC@orig@subsection\subsection
\let\MyTOC@orig@subsubsection\subsubsection
\let\MyTOC@orig@paragraph\paragraph

% SECTION
\renewcommand{\section}{%
  \@ifstar{\MyTOC@section@star}{\@dblarg\MyTOC@section@nostar}%
}
\def\MyTOC@section@star#1{%
  \MyTOC@orig@section*{#1}% sin entrada al .stc
}
\def\MyTOC@section@nostar[#1]#2{%
  \MyTOC@orig@section[{#1}]{#2}% comportamiento normal (escribe al .toc)
  % Escribimos también nuestra línea al .stc (texto con \numberline)
  \addcontentsline{stc}{section}{\protect\numberline{\thesection}#2}%
}

% SUBSECTION
\renewcommand{\subsection}{%
  \@ifstar{\MyTOC@subsection@star}{\@dblarg\MyTOC@subsection@nostar}%
}
\def\MyTOC@subsection@star#1{%
  \MyTOC@orig@subsection*{#1}%
}
\def\MyTOC@subsection@nostar[#1]#2{%
  \MyTOC@orig@subsection[{#1}]{#2}%
  \addcontentsline{stc}{subsection}{\protect\numberline{\thesubsection}#2}%
}

% SUBSUBSECTION
\renewcommand{\subsubsection}{%
  \@ifstar{\MyTOC@subsubsection@star}{\@dblarg\MyTOC@subsubsection@nostar}%
}
\def\MyTOC@subsubsection@star#1{%
  \MyTOC@orig@subsubsection*{#1}%
}
\def\MyTOC@subsubsection@nostar[#1]#2{%
  \MyTOC@orig@subsubsection[{#1}]{#2}%
  \addcontentsline{stc}{subsubsection}{\protect\numberline{\thesubsubsection}#2}%
}

% PARAGRAPH
\renewcommand{\paragraph}{%
  \@ifstar{\MyTOC@paragraph@star}{\@dblarg\MyTOC@paragraph@nostar}%
}
\def\MyTOC@paragraph@star#1{%
  \MyTOC@orig@paragraph*{#1}%
}
\def\MyTOC@paragraph@nostar[#1]#2{%
  \MyTOC@orig@paragraph[{#1}]{#2}%
  % Si no numeras los \paragraph, cambia \theparagraph por texto plano sin \numberline
  \addcontentsline{stc}{paragraph}{\protect\numberline{\theparagraph}#2}%
}

% --- Impresor del índice propio (lee el .stc) ---
\newcommand{\MyTOC}{%
  \par\bigskip
  {\centering\bfseries\Large Índice de contenido\par}%
  \medskip
  \begingroup
    \parindent=0pt\parskip=2pt
    \@starttoc{stc}%
  \endgroup
}

\makeatother
% ===================== fin ToC propio =====================


%%%%%%%%%%%%%%


% --- Datos del tema ---
\title{Tema 3.A.33 -- Teorías de la demanda de consumo corriente.\\
\large Ciclo vital y renta permanente. La demanda de bienes de consumo duradero. Evidencia empírica e implicaciones de política económica}
\author{Marcelo Ortega}
\date{Marzo 2026}
\newcommand{\TemaCode}{3A33}

\oldtitle{Tema 3.A.32. Teorías de la demanda de consumo. Implicaciones de política económica}
\changerationale{
    La demanda de consumo corriente difiere sustancialmente de la de bienes duraderos, tanto en la forma de modelizarla como en sus propiedades empíricas. Hasta ahora, la aproximación a este tema se centraba únicamente en la primera. En una crisis como la de la pandemia de COVID-19, los posibles fenómenos de "demanda embalsada" no pueden entenderse sin comprender primero cómo funciona la demanda de bienes duraderos.
}

\begin{document}


\maketitle
\changebox

\newpage

% --- Página de resumen y modificaciones ---
\puntosclave{ %% Escribir puntos clave Apuntes CECO
    Habría que añadir la demanda de bienes de consumo duradero y posiblemente más evidencia empírica. Además, sería necesario quitar peso a los modelos estáticos, que serían mencionados de pasada. También sería interesante ver si el consumo de bienes duraderos se podría incluir como salida a la contradicción entre las predicciones de Hall y la evidencia empírica.

    El enfoque keynesiano puede o bien incluirse en la exposición de manera breve a modo de antecedente (como se presenta en este tema) o bien obviarse y conocerse de cara a las preguntas del tribunal (como se presenta en el tema 3A34).
    }
\vspace{1cm}

\modificaciones{
    
    }

\newpage
\MyTOC
\newpage

\mylistofequations
\clearpage

\MyListOfFigures
\clearpage


\pagenumbering{arabic}
\setcounter{page}{1}


\section{Introducción}\label{intro}

En \textit{The Theory of Interest}, IRVING FISHER basa su explicación del ahorro personal en cinco características: previsión, autocontrol, hábitos, expectativa de vida y amor por la descendencia.


\subsection{Contextualización}\label{context}


\subsubsection{Relevancia}\label{importance}


\subsubsection{Historia}\label{history}



\subsubsection{Problemática}\label{problem} 




\subsection{Estructura de la exposición}\label{structure}
En primer lugar, analizaremos el consumo bajo certidumbre desde el modelo de Keynes hasta las teorías de la renta permanente de Friedman o la teoría de ciclo vital de Ando y Modigliani, los cuales se centran en las decisiones de consumo en un contexto dinámico. 

En segundo lugar, nos centraremos en un contexto con incertidumbre a partir del artículo pionero de Hall donde las expectativas de los agentes juegan un papel crucial. Por último, el análisis de la demanda de consumo va más allá, incluyéndose otros desarrollos como el consumo de bienes duraderos o el estancamiento secular.




\clearpage
\section{Demanda de consumo dinámica: Aproximacines ortodoxas}
\puntosclave{
    KEYNES introdujo la función de consumo como la relación entre el consumo y la renta. 
    
    No obstante, las irregularidades empíricas asociadas a su teoría llevaron a desarrollos posteriores como la aportación de DUESENBERRY quien afirmó que la satisfacción que obtiene un individuo del consumo no depende sólo de su \textbf{nivel absoluto} de renta, sino también de los niveles del entorno social con los que entra en contacto, es decir, de su \textbf{renta relativa}. 
    
    En un contexto dinámico se derivan las teorías de la renta permanente (FRIEDMAN) y la hipótesis del ciclo vital (ANDO y MODIGLIANI). Por una parte, FRIEDMAN se pregunta cuál es el nivel de consumo que un agente podría mantener permanentemente en el tiempo sin que se altere la riqueza del individuo. Por otra parte, Ando y Modigliani destacaron que la renta varía sistemáticamente a lo largo de la vida de los individuos y que el ahorro permite a los consumidores trasladar renta de las épocas de la vida en que ésta es alta a las épocas en que es baja.
}

A principios de los años 1950s se desarrollaron los dos modelos dominantes de consumo: la hipótesis de la renta permanente y la hipótesis del ciclo de vida. Aunque estos modelos de consumo inicialmente se consideraron como modelos antagónicos, a día de hoy se consideran complementarios. 

La idea central en torno a ambos es que los individuos no toman sus decisiones de consumo basándose únicamente en su renta actual, sino también en sus expectativas de ingresos futuros (perspectiva forward looking). Por ello, desde la década de 1950, surgen nuevas teorías basadas en un análisis intertemporal. 

Para ello, se añaden fuertes fundamentos microeconómicos a la decisión de consumo: el consumidor maximizará su utilidad distribuyendo el flujo de ganancias esperadas a lo largo de su vida en un esquema temporal de consumo, es decir, con perspectiva de \textit{forward-looking}.\footnote{%
  Se podría argumentar que KEYNES microfundamenta su formulación de la demanda de consumo mediante la ley psicológica fundamental y DUESENBERRY hace lo propio con los efectos demostración y trinquete. Por lo tanto, lo estrictamente correcto sería decir que a diferencia de los modelos previos, estos modelos se basan en programas de optimización de los agentes económicos.
}

\subsection{Hipótesis de la Renta Permanente: FRIEDMAN}
La hipótesis de la renta permanente hacía hincapié en las variaciones \textbf{estocásticas} de la renta a lo largo del tiempo. 

Así, el modelo propuesto por FRIEDMAN descompone al consumo en dos partes: permanente ($C_t^P$) y temporario ($C_t^T$); como consecuencia de existir dos tipos de renta: permanente ($Y_t^P$) y temporario ($Y_t^T$). ¿Cómo se relacionan entre ellos? 

Para FRIEDMAN, la renta permanente es una porción de la riqueza total del individuo: $Y_t^P=r\cdot W_t$. Puede entenderse como el flujo de renta que puede ser consumido en un periodo de tiempo \textbf{sin afectar al valor de la fuente de la cual proviene}. Por ello, se concibe (en su formulación más básica), como el rendimiento obtenido en el periodo $t$ por la riqueza acumulada por el individuo hasta ese momento ($W_t$), compuesto por activos reales, activos financieros (dinero, bonos), y riqueza humana siendo. No es de extrañar, por tanto, que la estimación empírica de la riqueza se convirtiera en una cuestión central en la especificación de la función de consumo de Friedman; así como también en su principal debilidad. 

Por otro lado, FRIEDMAN concibe la renta transitoria como una variable compuesta por dos elementos: (i) shocks transitorios que afectan a la renta; y, (ii) incrementos contingentes de renta que serán integrados como permanentes. En ambos casos, se formulan como shocks estocásticos.

Respecto al consumo, FRIEDMAN postula la existencia de un consumo permanente proporcional a la renta permanente. Este coeficiente proporcional, por ejemplo $k$, depende de factores tales como el tipo de interés, los gustos y elementos demográficos. Y, por otro lado, el gasto de un periodo cualquiera ($C_t$), se descompone en un elemtno permanente ($C_t^P$), que ya hemos definido, y un componente transitorio. ¿Por qué? Por ejemplo, este elemento transitorio puede racionalizarse como periodo de compras extraordinarias, gastos de enfermedad, etc. 

Con esta definición, podemos ver como FRIEDMAN enfatiza la noción de redistribución intertemporal del consumo: el agente \textit{siempre quiere} consumir prácticamente la misma porción, y son los sucesos contingentes lo que le separan de esta tendecia lineal. 

\subsubsection{Agregación: ¿cómo se obtiene el consumo agregado?}
Hemos definido el comportamiento individual de una persona cualquier. Pero, lo que FRIEDMAN pretende es dar respuesta a los patrones de consumo agregado, por tanto: ¿cómo podemos agregar estos patrones individuales de consumo? Para conseguir la agregación debemos añadir una serie de supuestos.


En primer lugar, supondremos que los shocks transitorios de consumo se cancelan entre sí: los individuos tendrán shocks positivos y otros negativos, aunque a nivel individual puedan ser muy importantes. 

En segundo lugar, supondremos que, a nivel agregado, las rentas transitorias positivas y negativas de los agentes no se cancelan en su totalidad. Es decir, existe un sesgo sistemático que empuja la renta de la economía hacia arriba.

En tercer, y último lugar, asumiremos que los componentes del consumo, transitorios y permanentes, no están correlacionados entre si.

\paragraph{Sección cruzada: Representación gráfica}
Gracias a estos supuestos FRIEDMAN es capaz de reconciliar su teoría con los datos de sección cruzada (KUZNETS y DUESENBERRY, [\authorfont{Ver Anexo \ref{anx:antecedentes_historicos}}]), que muestran unas relaciones de consumo-renta más altas en los tramos de renta más bajos. 

Es decir, como en los grupos de renta más bajos habrá una mayor proporción de individuos con rentas transitorias negativas que en los grupos de rentas altas, la media de las rentas transitorias será negativa en el primer grupo y positiva en el segundo, lo que conlleva una propensión al consumo por encima y por debajo de la media, respectivamente. Gráficamente:

\imagenfit{PIH_FRIEDMAN.png}{Consumo bajo la Hipótesis de la Renta Permanente}{\textit{Macroeconomía Avanzada I}. Argandoña, A.; Gámez, C. y Mochón, F. (1997)}

En el gráfico podemos ver la función de consumo de un agente representativo que, partiendo de un punto $D$ donde el y la renta transitorias son nulas, experimenta una serie de shocks. Por un lado, la curva de consumo originado por la renta permanente mantiene una pendiente $k$ estructural. Por otro lado, el componente de renta transitoria ($Y_t^T$) aplana la curva de consumo observado que observamos en el gráfico.

Supongamos que el agente tiene una renta transitoria positiva situándose en el punto $A$, de forma que, como el consumo depende principalmente del nivel de renta permanente, se observa que el incremento en el consumo es mucho menor que el incremento transitorio de la renta: el agente ahorra o invierte ese aumento transitorio de la renta manteniendo el mismo nivel de consumo. Asimismo, si el agente se encontrase en $B$, con un shock de renta transitorio negativo, la disminución de renta no implicaría una disminución del consumo en igual cuantía: el agente desahorro o pide prestado para mantener un nivel similar de consumo. Esto es precisamente lo que se conoce como \textbf{suavización del consumo}.

Además, nótese que el mismo gráfico permite de hecho explicar la evolución del consumo a largo plazo en relación al ciclo económico: en las fases alcistas la media de rentas transitorias es positiva ($\mathbb E[Y_t^T]>0$) mientras que lo contrario ocurre en periodos de recesión, pero la relación consumo-renta a largo plazo sería estable y vendría dado por la renta permanente. 

\subsubsection{Implicaciones de Política Económica}
De este modo, sólo en la medida en que estos componentes transitorios de renta se incorporen a la renta permanente (para lo que debe transcurrir cierto tiempo) afectarán al consumo.

Si bien Friedman llegó a estimar que una variación de la renta corriente total del periodo habría de persistir durante tres años para trasladarse totalmente a la renta permanente, no han podido ser totalmente resueltos los graves problemas de contrastación empírica de esta teoría, haciéndose necesario el estudio del impacto de las expectativas racionales sobre el consumo, tal y como se haría años más tarde.

En suma, de esto se desprende el poco poder que atribuía FRIEDMAN a la Política Fiscal, siendo preferida la Política Monetaria puesto que un aumento de la oferta monetaria tendría un efecto riqueza sobre el consumo a corto plazo, atendiendo a los supuestos previamente comentados.

\paragraph{Limitaciones: ¿qué queda fuera de esta aproximación?}
Sin embargo, esta teoría presenta una serie de inconvenientes: (i) los principales componentes de la riqueza no humana no tenían valoración de mercado; (ii) las estimaciones de riqueza no formaban parte de las cuentas de la renta nacional; (iii) a efectos de previsión, se requerían importantes ecuaciones para pronosticar la evolución futura de la riqueza; y, (iv) no resuelve completamente la relación entre la renta del periodo y el consumo durante el mismo.\footnote{%
  KOTLIKOFF y SUMMERS en 1981 presentaron estudios que indicaban cómo el ahorro no es tanto origen de consumo futuro como un medio de dejar herencia
} Además, FRIEDMAN asume que los agentes tienen certidumbre sobre las rentas futuras, pero la realidad es que los agentes toman las decisiones en un contexto de incertidumbre. 

No obstante, más allá de sentar las bases para teorías posteriores, la teoría de la renta permanente de FRIEDMAN, mantiene todavía cierta importancia en la actualidad. 

Así, por ejemplo, HUBBARD, SKINNER y ZELDES (1994, 1995) explican por qué la mayoría de los hogares conservan poca riqueza a partir de una hipótesis muy próxima al espíritu de la renta permanente. El elemento clave de su explicación, aparte de la optimización intertemporal, es el hecho de que los programas de asistencia social ofrecen un seguro frente a niveles muy bajos de consumo. Estos programas desincentivan el ahorro de aquellos hogares con una probabilidad significativa de tener que recurrir a ellos de dos maneras: (i) porque proporcionan una forma de seguro frente a la posibilidad de que se  cumplan las expectativas de ingresos menos favorables; y, (ii) porque imponen un elevado tipo impositivo implícito sobre la tenencia de activos. 


\subsection{Teoría del Ciclo Vital: ANDO y MODIGLIANI}
Si, como hemos visto, FRIEDMAN basa las fluctuaciones del consumo y renta en variaciones estocásticas, la Hipótesis del Ciclo Vital se centra en las variaciones \textbf{predecibles} de la renta a lo largo del ciclo vital para explicar las fluctuaciones en el consumo.

Esta hipótesis, debida a ANDO y MODIGLIANI (1963) y ANDO y BRUMBERG (1954), utiliza para su desarrollo el modelo de FISHER, sobre decisiones intertemporales, para determinar el consumo de una persona en función de la renta/riqueza que prevee percibir a lo largo de toda su vida. 

\textcolor{magenta}{Explicar con palabras la derivación que había antes aquí y que se ha pasado a los anexos}

Como vemos, esta formulación no es más que la representación analítica del siguiente patrón de consumo representado gráficamente:\footnote{%
    Para simiplificar la ecuación eliminando el sumatorio, i.e. $H_{i,t}\to \Phi_{i,t}y_{i,t}$, seguimos los siguientes pasos. En primer lugar, cambiamos los indices del sumatorio $j=t$ por $j=t+s$: $W_{i,t}=A_{i,t-1}+\sum_{s=0}^{T-t}\frac{\mathbb E[y_{i,t+s}]}{(1+r)^s}$. En segundo lugar, sustituimos por las HEA $\mathbb E[y_{i,t+s}]=\phi_s\,y_{i,t}$: $W_{i,t}=A_{i,t-1}+\sum_{s=0}^{T-t}\frac{\phi_s\,y_{i,t}}{(1+r)^s}$. Sacamos $y_{i,t}$ fuera del sumatorio: $W_{i,t}=A_{i,t-1}+y_{i,t}\cdot \sum_{s=0}^{T-t}\frac{\phi_s}{(1+r)^s}$. Y, por último, definimos: $\Phi_{i,t}=\sum_{s=0}^{T-t}\frac{\phi_s}{(1+r)^s}$
}

\imagenfit{HCV_MODIGLIANI.png}{Consumo e ingresos a lo largo del ciclo vital}{\textit{Macroeconomics}. MANKIW, N.G. (2019)}

El postulado principal sobre el que se basa este modelo es que los consumidores pretenden poseer un nivel constante de calidad de vida a lo largo de ésta. Sin embargo, la renta varía sistemáticamente a lo largo de la vida de los individuos y es el ahorro (y endeudamiento) el que permite a los consumidores trasladar rentas intertemporalmente.

En los primeros años de vida, los ingresos laborales suelen ser bajos en relación con los años de trabajo posteriores, mientras que los ingresos suelen alcanzar un máximo en la última parte de la vida laboral para descender en la jubilación. Por tanto, como los consumidores desean suavizar el consumo preferirán endeudarse en los primeros años, reembolsar esos préstamos y acumular riqueza durante los años de ingresos altos, y luego desahorrar durante la jubilación. 

En otras palabras, con el fin de mantener constante el nivel de consumo durante toda la vida el consumidor ahorra durante su vida activa para desahorrar en la vida pasiva. 

\subsubsectio{Agregación: ¿cómo ...?}
Igual que antes, debemos agregar los patrones de consumo individual de los agentes, de acuerdo con esta Teoría, con el fin de obtener el consumo agregado. 

\textcolor{magenta}{Para ello... Imponiendo una serie de parámetros adicionales respecto a las distribuciones de determinadas variables (propensiones medias al consumo, distribuciones de riqueza y renta y estructura demográfica), podemos obtener la siguiente función agregada de consumo como: (continuar, aunque esta agregación no sea literal, que se pueda interpretar conforme a la tésis de los autores)}


La función agregada de consumo que obtendríamos se puede representar gráficamente como:

\textbf{¿INCLUIR GRÁFICOS?}

En primer lugar, vemos la primera etapa de la vida, con una renta muy reducida y un consumo elevado, en términos relativos (elevada propensión media al consumo, curva $C_1$). A medio y largo plazo el individuo va acercándose a la segunda etapa de su ciclo vital, por lo que su renta aumenta, pasando de $Y_1$ a $Y_2$. Al aumentar su renta, aumenta su riqueza pasando de $A_{t-1}$ a $A_{t-1}'$, lo que hace que la curva $C_1$ se desplace hacia arriba, $C_2$. Uniendo los puntos $A$, $B$ y $C$, se obtendría la curva de consumo a largo plazo, donde la propensión media al consumo coincide con la propensión marginal al consumo. Por tanto, cuanto menores sean los niveles de renta del grupo de población, mayor será la propensión media al consumo, ya que se trataría de los grupos poblacionales de jóvenes y ancianos; y viceversa.

\textcolor{magenta}{No entiendo muy bien si lo que viene a continuación es algo previo a la agregación o no. Decide donde crees que es más adecuado ponerlo. Gracias.}

Como podemos intuir, la Hipótesis del Ciclo Vital permite tambié explicar la tendencia al descenso de la propensión media al consumo ($\frac{C}{Y}\equiv PMe_C$). 

Tanto desde un punto de vista intergeneracional como desde un enfoque \textcolor{red}{intercalases}, se observa que las familias en los estratos de influjos de renta más bajos presentan una elevada propensión media al consumo (jóvenes y ancianos), mientras que en los estratos de influjos de renta más elevados encontraremos una propensión media al consumo más baja (edad media). Gráficamente:

\imagenfit{PMeC_HCV.png}{Estudio de corte trasversal}{\textit{Macroeconomía Avanzada I}. ARGANDOÑA, A.; GÁMEZ, C. y MOCHÓN, F. (1997)}

\textcolor{magenta}{¿Por qué? ... (continuar, explicar esto conforme a la Teoría y por qué se puede observar, no entiendo muy bien la relación..)}


\subsubsection{Implicaciones de Política Económica}
En definitiva, la Teoría del Ciclo Vital concuerda con las observaciones empíricas de corto plazo en lo que respecta a la propensión marginal a consumir, menor que la propensión media al consumo. También fundamenta el comportamiento a largo plazo del consumo a través de la relación entre consumo y renta y el hecho de que la propensión media al consumo a largo plazo sea constante. Además, incluye explícitamente los activos como una variable explicativa de la función de consumo. 

Las conclusiones de Política Económica son, semejantes a las derivadas del modelo de renta permanente de FRIEDMAN, estando en este caso más justificada la introducción de la riqueza.\footnote{%
    \href{https://www.nobelprize.org/uploads/2018/06/modigliani-lecture.pdf?utm_source=chatgpt.com}{MODIGLIANI, Discurso del Premio Nobel (2018)}
}

Por un lado, a corto plazo: (i) el hecho de que la riqueza tenga un papel importante en la función de consumo de corto plazo implica que la \textbf{política monetaria puede afectar a la demanda agregada} no solo a través del canal tradicional de la inversión, sino también \textbf{a través del valor de mercado de los activos y del consumo}; y, (ii) los intentos de frenar (o estimular) la demanda mediante \textbf{impuestos transitorios sobre la renta} (o devoluciones fiscales) pueden esperarse que tengan \textbf{efectos pequeños sobre el consumo} y que reduzcan (o aumenten) el ahorro, porque el consumo depende de los recursos a lo largo de la vida, los cuales apenas se ven afectados por un cambio impositivo transitorio (resultado respaldado empíricamente). 

Por otro lado, algunas proposiciones que se derivan del modelo son: (i) un \textbf{impuesto progresivo sobre el consumo es más equitativo} que uno sobre la renta corriente porque se ajusta más a la renta permanente, independientemente de sus efectos de incentivo sobre el ahorro; y, (ii) el gasto financiado mediante déficit tiende a ser pagado por generaciones futuras, mientras que el financiado mediante impuestos lo paga la generación actual.\footnote{%
    Esta conclusión se basa en la proposición de que el ahorro privado, al estar determinado por consideraciones del ciclo vital, debería ser (casi) independiente de la posición presupuestaria del gobierno (MODIGLIANI y STERLING [1985]), y por lo tanto la riqueza privada debería ser independiente de la deuda pública (MODIGLIANI, 1984).
}.

De esto último se deduce que la deuda pública tiende a desplazar (crowding out) una cantidad equivalente de capital privado, con un coste social igual al rendimiento del capital perdido (que también es aproximadamente igual al pago de intereses de la deuda pública). Esta conclusión contrasta claramente con la defendida por la llamada Proposición de Equivalencia Ricardiana (BARRO, 1974), que sostiene que siempre que el gobierno incurre en déficit, el sector privado aumentará su ahorro para compensar el efecto desfavorable del déficit sobre las generaciones futuras. Por supuesto, en la medida en que el déficit público se utilice para financiar inversiones productivas, las generaciones futuras también recibirán los beneficios del gasto, y permitir que lo paguen mediante financiación con déficit puede ser compatible con la equidad intergeneracional.

Además, en una economía abierta, el efecto de desplazamiento de la inversión puede verse atenuado por la entrada de capital extranjero, atraído por el mayor tipo de interés que resulta de la menor disponibilidad de fondos invertibles. Sin embargo, la carga para las generaciones futuras permanece aproximadamente inalterada debido a los intereses que deberán pagarse sobre la deuda externa.

Por otro lado, si existe capacidad ociosa en la economía, el gasto público financiado mediante deuda puede no desplazar la inversión (al menos si va acompañado de una política monetaria acomodaticia) y puede, en cambio, elevar la renta y el ahorro. En este caso, el déficit resulta beneficioso, como sostenían los primeros keynesianos; sin embargo, la deuda tendrá un efecto de desplazamiento una vez que la economía vuelva al pleno empleo. La HCV sugiere que, para evitar este resultado, existe un buen argumento a favor de un presupuesto equilibrado cíclicamente.

\textcolor{magenta}{LIGAR AQUÍ LA EXPLICACIÓN DE ANTES (forma parte del mismo estudio): HUBBARD, SKINNER y ZELDES (1994, 1995)} En cuanto a los hogares cuyas perspectivas de renta son lo bastante favorables como para descartar la posibilidad de tener que recurrir a los programas de asistencia, el consumo viene determinado por la optimización intertemporal, y por tanto exhiben las pautas de ahorro previstas por la hipótesis del ciclo vital. 

De ahí que HUBBARD, SKINNER y ZELDES afirmen que los diferentes modelos de acumulación de riqueza de ricos y pobres pueden ser explicados sin necesidad de acudir a diferencias en sus preferencias.

\paragraph{Limitaciones: ¿qué está pasando?}
Una de las características más simples y elegantes del modelo del ciclo vital es la forma en que maneja la variabilidad de la renta: en cada período el consumidor debería consumir el valor anualizado de su riqueza esperada. 

Téngase en cuenta que esta afirmación es válida tanto si la variabilidad de la renta es determinística como estocástica. Por ello es hoy en día considerado como complementario a la Hipótesis de la Renta Permanente, que consideraba únicamente variaciones estocásticas. 

No obstante, nótese que, asumiendo ésto estamos suponiendo que los consumidores tienen algún tipo de \textbf{expectativas racionales}.\footnote{%
  \textcolor{magenta}{¿Por qué?... Explicar aquí por qué con eso asumimos lo otro (relación causal)}
} Por tanto, los aumentos permanentes de la renta generan respuestas mucho mayores del consumo que los aumentos transitorios, porque implican incrementos mayores de la riqueza. 

Bajo estos supuestos (explícitos e implícitos) y en un contexto de plena certidumbre, la demanda de consumo dinámica del consumidor se resuelve de una sola vez: eligiendo las sucesiones de consumo y estando únicamente sujetas a variaciones de la naturaleza expresada. Esta es precisamente la aproximación teórica vista hasta ahora. 

Sin embargo, en un contexto de incertidumbre como el que caracteriza nuestra realidad económica, las decisiones del consumidor se resuelven en cada período en base a la información razonablemente predecible: en cada período se obtiene información nueva que afecta a las expectativas sobre valores futuros.


\subsection{Paseo aleatorio: HALL}
Precisamente esta cuestión fue aborada poco después por el economista HALL (1978).

HALL buscó desarrollar el primer modelo en incorporar la HER dentro de las Teorías de demanda de consumo corriente. De esta forma, se resolvería la cuestión que dividía por aquella época ambas aproximaciones y, al mismo tiempo, se aproximaba ambas teorías sustancialmente más a la realidad, incorporando la expectativas futuras mediante HER.

Partiendo de un problema de optimización intertemporal clásico podemos derivar la condición de primer orden como la igualdad entre utilidades marginales de dos periodos correlativos, el corriente y el siguiente esperado. 

\textcolor{magenta}{Hay que desarrollar aquí en palabras lo que se mostraba antes en fórmulas (y se ha relegado al anexo). El objetivo es primar la intución económica más que el desarrollo matemático. (elaborar)}

En definitiva, el Modelo de HALL permite reconciliar la Hipótesis de la Renta Permanente y la Hipótesis del Ciclo Vital, unificándolas bajo un mismo paraguas e incorporando la HER. 

\textcolor{magenta}{(Hay que seguir aquí, manteniendo lo que se quiere decir, pero explicandolo mejor porque no queda nada claro (yo no lo entiendo)). Esto se produce porque HALL y MISHKIN separan la renta de los hogares en tres componentes: (i) un \textbf{componente determinístico} $y_{Dt}$, que aumenta con la edad hasta justo antes de la jubilación; (ii) un \textbf{componente estocástico} $y_{Lt}$, que fluctúa a medida que cambian las perspectivas de ingresos a lo largo de la vida y que se especifica como un paseo aleatorio (random walk); y, (iii) un \textbf{componente estocástico estacionario} $y_{st}$, que fluctúa debido a influencias transitorias y se describe mediante un proceso de media móvil.}

\subsubsection{Implicaciones de Política Económica}
En síntesis, el modelo de HALL concluye que: (i) para predecir el consumo futuro ($t+1$) solo necesitamos el consumo presente ($t$); (ii) los valores retardados de cualquier otra variable que en principio pudiesen ser relevantes para predecir el consumo futuro ($t+1$) no tienen valor explicativo adicional una vez que consideramos el consumo presente ($t$); (iii) el término de error refleja la nueva información sobre rentas futuras, información que no estaba disponible en el presente ($t$).

Por tanto, ¿cuáles son las implicaciones de Política Económica?

En primer lugar, desde la óptica del Modelo de HALL, la Política Económica modificará las decisiones de consumo a través de su impacto sobre la renta permanente, si y solo sí estos efectos sobre la renta son inesperados. Si, por el contrario, los cambios en la renta permanente son esperados no habrá ninguna variación en la estimación de la renta permanente, y ninguna variación en el consumo del agente. 

Por otro lado, el consumo no mostrará un comportamiento predecible a lo largo del ciclo. Cuando el output cae inesperadamente, llevará a una caída inesperada de la renta permanente y el consumo caerá. Si luego el output presenta una tendencia alcista, al no haber cambio inesperado de la renta permanente el consumo seguirá estable, por lo que no se recuperará. 

Por último, cualquier medida de Política Económica que conduzca a un aumento inesperado de la renta permanente no solo tendrá efectos sobre el consumo presente, sino sobre toda la senda futura de consumo, pues no presenta reversión a la media (no hay media, de hecho). 

\paragraph{Limitaciones: Contrastación Empírica}
La cualidades más importante del Modelo de HALL era, precisamente, la derivación de una regresión fácilmente contrastable con los datos empíricos. Por ello, los propios autores testaron el modelo contrastándolo con observaciones en uno de los primeros trabajos que se realizaron. 

Los autores estaban interesados en el parámetro $\beta_1\in[0,1]$, que refleja la propensión marginal a consumir a partir de la renta transitoria ($y_{st}$). 

Como es normal, debido a la HER e intertemporalidad del consumo, el modelo predice que el parámetro ($\beta_1$) será igual al valor anualizado de un dólar de renta transitoria.\footnote{%
  \textbf{NOTA.-} Esto es fundamental a la hora de entender bien el modelo. En HALL-MISHKIN, un dolar de renta puramente transitoria no debe traducirse en un aumento grande del consumo corriente, porque no altera (apenas) la riqueza vitalicia: solo añade un dolar al stock de riqueza en ese instante. Por tanto, el consumidor óptimo no se gasta ese dólar entero hoy, sino solo la \textbf{anualidad que le corresponde}. 

  Analíticamente, podemos verlo de manera sencilla. Denotando $\gamma_t$ el valor de una unidad de riqueza, descontado al presente (anualizado), entonces el consumo satisface: 

  $$c_t=\gamma_t(H_t+A_t), \qquad \gamma_t=\left(\sum_{\tau=0}^{T-t}(1+r)^{-\tau}\right)^{-1},$$

  Donde $H_t$ es riqueza humana (renta) y $A_t$ riqueza no humana (activos). Por eso, la propension marginal a consumir a partir de un dolar transitorio es precisamente $\gamma_t$: un dolar hoy, que no se repite en el futuro, solo eleva el consumo en la cantidad imputable a ese año. 
      
  Por tanto, con horizonte infinito, $\gamma_t\approx r$. Pero, con horizonte finito, \textbf{$\gamma_t$ depende de $r$ y de los años de vida restantes}. Eso es lo que quiere decir que $\beta_1$ deberia ser igual al valor de un dolar de renta transitoria, descontado al presente (anualizado): $\beta_1=\gamma_t$, para ser estructuralmente coherente con el modelo.
} La intuicion economica es muy limpia: (a) si la sorpresa de renta es permanente (\textbf{cambios en renta determinística o tendencia}, $\beta_0$.), afecta a toda la senda esperada de ingresos futuros, aumenta mucho la riqueza permanente, y  el consumo debe aumentar bastante; pero, (b) Ss la sorpresa es transitoria (\textbf{cambios en renta transitoria}, $\beta_1$), la riqueza total apenas cambia en un dólar, y el consumo solo aumenta en la anualidad de ese dólar. 

Dicho de manera clara: para HALL, el consumo solo debe moverse cuando llega nueva informacion relevante sobre la renta permanente. Por eso, más allá del consumo tendencial, el consumo efectivo de un periodo no debe ser predecible con informacion pasada como la renta rezagada
    
Como deberíamos haber notado, esta asimetria tan pronunciada entre shocks permanentes y transitorios es, precisamente, el nucleo común de la Hipotesis de la Renta Permanente, la Hipotesis del Ciclo Vital y la version de HALL con HER. 

Ahora bien, ¿validan las observaciones dicha predicción? HALL y MISHKIN estimaron para el consumo de alimentos utilizando el Panel Study of Income Dynamics para el período 1969–1975, el valor de $\beta_1$, y este resultó ser $0,29$. 

Interpretado literalmente, dice que de cada dolar de renta transitoria el hogar gasta unos $29$ centavos en ese año en consumo (de alimentos). Para reconciliar eso con la regla óptima, haria falta una anualidad $\gamma_t$ anormalmente alta, es decir, o bien \textbf{tipos de interés muy elevados o bien un horizonte restante muy corto}.\footnote{%
    SHREFFIN y THALER consideran que esto precisamente constituye una reconfirmación de las estimaciones anteriores de FRIEDMAN y HOLBROOK, que situaban las tasas de descuento en el intervalo $0,33–0,50$. \textit{Es notable que obtengan este resultado a pesar de utilizar el consumo de alimentos como variable dependiente. El consumo de alimentos parecería ser menos volátil que otros componentes del consumo. (...)}
} 

Como ninguna de las dos alternativas parece natural para la muestra, el resultado sugiere que el modelo puro no basta. HALL y MISHKIN extraen de ahi una conclusion muy concreta: una fraccion del consumo parece fijarse como proporcion de la renta corriente, no por la regla intertemporal optima completa. 

Los autores, y otros posteriores, definieron esta \textit{anomalía} como la \textbf{hipersensibilidad del consumo}: la violacion flagrante de la suavización del consumo debido a shocks transitorios de renta.

La hipersensibilidad del consumo significa que éste reacciona demasiado a la renta corriente o a la renta transitoria, mas de lo que permite el modelo de suavizacion intertemporal. Si $\beta_1$ sale muy por encima de $\gamma_t$, el consumidor esta comportandose como si una parte excesiva del ingreso transitorio se consumiera de inmediato.\footnote{%
    HALL y MISHKIN encuentran justamente eso: el consumo responde mucho mas a la renta permanente que a la transitoria, pero la respuesta a la transitoria sigue siendo notable, y solo encaja con el modelo si los tipos de interes son muy altos; ademas, encuentran que la covariacion observada es compatible con conducta ciclo vital/renta permanente para alrededor del 80\% del consumo, y con simple proporcionalidad a la renta corriente para el 20\% restante.
}
    
Este es solo uno de los muchos estudios empíricos que contrasto este fantasma que acechaba las Teorías de la demanda de consumo corriente basadas en la optimización racional. Para resolverlo, muchas teorías fueron presentadas con el fin de reconciliarlo, entre otras, el ahorro precautorio, las restricciones de liquidez o la demanda/stock de bienes duraderos.\footnote{%
  Los propios autores, tras un examen más detenido, encontraron que el $20\%$ del consumo total (de alimentos) no es explicado por el modelo del ciclo vital, y en consecuencia plantearon la hipótesis de que se fija como una fraccion de la renta corriente en lugar de seguir la regla óptima mas compleja. Esto les llevo a señalar que no eran capaces de distinguir este síntoma de incapacidad (o falta de voluntad) para pedir prestado y prestar de aquel tipo de comportamiento característico de consumidores que simplemente afrontan tipos de interés elevados. 

  Esta observación final de HALL y MISKIN sobre incapacidad o falta de voluntad para pedir prestado y prestar va exactamente por ahí. Si un hogar no puede endeudarse, o no quiere hacerlo, entonces un ingreso transitorio positivo eleva el consumo mas de lo que predice el modelo frictionless, porque ese ingreso relaja una restriccion de liquidez actual. En ese caso, una $\beta_1$ alta no implica necesariamente irracionalidad; puede ser el reflejo de mercados de credito imperfectos o de aversion al endeudamiento. Ellos mismos dicen que no pueden distinguir bien entre ambas hipótesis con sus datos. 
}

En definitiva, podría argumentarse que este trabajo de Hall se sustenta sobre supuestos muy restrictivos. Por ejemplo, las conclusiones alcanzadas no se mantendrían si la utilidad marginal no fuese lineal. Por ello, se trató de verificar empíricamente hasta qué punto estos corolarios podían defenderse con la evidencia. 


Lo que se trató de verificar fundamentalmente es si los cambios predecibles en la renta producen cambios en el consumo o no. En caso de que el consumo sí que responda a cambios predecibles en la renta se estaría produciendo lo que se conoce como hipersensibilidad del consumo. Destacamos fundamentalmente dos trabajos


\subsection{Demanda de bienes de consumo duraderos}
Como hemos visto, este trabajo extrajo sus conclusiones a raíz del análisis de datos sobre alimentos: resulta difícil encontrar un bien de consumo más perecedero. 

Por ello resulta interesante presentar una extensión que pretende resolver este fenómeno a la vez que proporcionar una Teoría de la demanda de consumo de bienes duraderos. 

\textcolor{magenta}{(Reformular para explicar en palabras cuál es el programa de optimización que sucede con bienes durables, y qué es lo que cambia)La diferencia clave entre los bienes no duraderos y los duraderos es que los primeros son una variable flujo, mientras que los segundos son una variable stock que proporciona un flujo de servicios (utilidad) en cada periodo de vigencia. De esta forma, deben ser incluidos de manera explícita en el problema de maximización del consumidor. (...)}

Podemos representar gráficamente la composición del consumo del agente con la siguiente:

\textbf{¿SE PUEDE INCLUIR UN GRÁFICO?}

Esta nueva formulación tiene tres implicaciones principales.

En primer lugar, la \textbf{desvinculación consumo-gasto}. El consumo relevante no coincide con el gasto cuando hay bienes duraderos. Es decir, el consumo en no duraderos se eligen como flujo periodo a periodo, lo que corresode con el gasto. En cambio, los duraderos se eligen como stock y por tanto implican una decision de cartera intertemporal, desvinculando el \textit{consumo} del \textit{gasto}. Desde la óptica del investigador: el gasto en duraderos será mucho más volátil y más sensible a tipos de interés, renta esperada y precios relativos que el gasto en no duraderos. No porque el consumidor suavice peor, sino porque el \textbf{gasto en duraderos contiene un componente de inversión}

En segundo lugar, la relación de \textbf{composición óptima de consumo}. La composición óptima entre no duraderos y servicios del stock durable depende del coste de uso del durable ($p_t^D$), no solo de la renta. No obstante, los cambios sobre la renta permanente también tienen proyección sobre el modelo. En no duraderos, el efecto es claro: mayor riqueza permanente implica mayor flujo de consumo desde hoy. En duraderos, el hogar deseará un stock objetivo mayor. Para acercarse a ese nuevo stock objetivo, aumenta la compra de duraderos hoy; pero como el stock se acumula y se deprecia lentamente, el consumo será gradual incluso si el gasto pega un salto (e $t$). Por eso el gasto en duraderos puede sobrerreaccionar al principio (hipersensibilidad).

En tercer lugar, el \textbf{ajuste empírico}. Los shocks de renta o de tipos de interés pueden generar respuestas muy distintas en no duraderos y en gasto en duraderos, porque este ultimo es ajuste de stock. Tal y como se ha anticipado en las dos implicaciones anteriores. En particular, el mejor ajuste empírico se origina en: (i) las \textbf{perturbaciones en renta}, pues parte de la hipersensibilidad se origina en las observaciones de compras (gasto) de bienes duraderos, que no son un flujo puro de consumo sino un ajuste de stock; y, (ii) los \textbf{retardos}, pues incorporando duraderos aparecen de forma natural retardos de ajuste. Por tanto, no hace falta suponer miopía para explicar por que el gasto no se mueve uno a uno con la nueva riqueza permanente. Basta con reconocer que el stock durable heredado del pasado condiciona la elección presente.

De esta forma, trabajar con no duraderos y servicios permitiría evitar que la dinámica de los bienes duraderos contamine la contrastación del modelo básico. Una renta transitoria puede provocar una reacción fuerte del gasto en duraderos aunque el modelo intertemporal sea correcto. 

Se podría decir que la incorporación de bienes duraderos interpreta que si el parámetro de contrastación de HALL ($\beta_1$) es elevado, demuestra que la variable observada contiene un componente que no sigue la formulación básica para no duraderos. Lo que no contradeciría el fenómeno de la suavizacion del consumo; sino que explicaría por qué el gasto observado es \textit{hipersensible}, mezclando consumo y acumulación de stock durable. 

De hecho, los hogares pueden reducir temporalmente las compras de bienes duraderos cuando sus ingresos disminuyen, sin que se produzca una gran reducción de la utilidad a corto plazo. Cuando sus ingresos se recuperan, pueden reanudar sus compras para compensar las pérdidas anteriores en su stock de bienes duraderos.

\subsubsection{Implicaciones de Política Económica}
A partir de los resultados de diversos estudios empíricos centrados en la respuesta relativa de los bienes de consumo duradero y no duradero ante cambios en la renta, se observa que los bienes duraderos presentan una mayor volatilidad que los no duraderos. 

MANKIW señala que la desviación estándar de la tasa de crecimiento de los bienes no duraderos es sólo del $1,4\%$, mientras que la desviación estándar de los bienes duraderos es del $8,5\%$. Por tanto, los bienes duraderos desempeñan un papel fundamental en el ciclo económico. 

Además, los bienes duraderos también tienen una mayor tendencia a financiarse mediante crédito, debido a su elevado valor y a su mayor vida útil. 

Por otra parte, debido a su mayor vida útil, algunos bienes duraderos pueden utilizarse como garantía, lo que facilita su financiación mediante el crédito. Sin embargo, en tanto las normas de crédito suelen ser más estrictas durante las recesiones económicas y más laxas durante las expansiones, esto también aumenta la sensibilidad de los bienes duraderos al ciclo económico.

Una forma de explicar la excesiva sensibilidad del consumo respecto a la renta corriente encontrada en numerosos trabajos empíricos se basa en la no separabilidad de las preferencias de los individuos, es decir, el consumidor deriva utilidad tanto de los bienes no duraderos de consumo como de los servicios que generan los duraderos de forma que estas decisiones no son independientes.\footnote{%
    En consecuencia, la teoría del ciclo vital y la teoría de la renta permanente deberían incorporar separadamente ambos tipos de bienes, determinándose, en este caso,  conjunta, y no separadamente, las sendas óptimas del consumo de cada uno de ellas.
} 

La implicación que deduce BERNANKE (1985), a partir de una función de utilidad no separable en bienes duraderos y no duraderos de consumo y considerando costes de ajuste del stock de bienes duraderos, es que el consumo de bienes no duraderos no sigue un paseo aleatorio o un proceso autorregresivo de primer orden. 

Varios estudios reconocen que la decisión de adquirir bienes duraderos puede constituir un mecanismo para preservar la suavidad del gasto en bienes y servicios no duraderos. A la hora de decidir el importe de la compra y el momento en que se realiza, el consumidor puede adelantar o posponer el gasto en función de determinados factores. Es esta posibilidad de decidir cuándo van a realizar el gasto la que permite a los consumidores utilizar la adquisición de bienes duraderos como mecanismo de absorción de perturbaciones.


\paragraph{Limitaciones: ¿valida la contrastación?}

\textcolor{magenta}{Incluir aquí una serie de limitaciones que tengan la Teoría de demanda de bienes duraderos. Más que nada para seguir con la estructura de los apartados anteriores...(continuar)}



























\clearpage
\section{Demanda de consumo dinámica: Aproximaciones heterodoxas}
\puntosclave{

}

La Hipótesisdel Civlo Vital y la Hipótesis afín de la Renta Permanente de FRIEDMAN (1957), así como la incorporación de HER propuesta por HALL, son ejemplos clásicos de teorización económica. Todas estas formulaciones introducen, explíctia o implíctamente, supuestos simplificadores con el fin de caracterizar un problema de optimización bien definido, que posteriormente se resuelve. La solución a ese problema de optimización constituye el núcleo de la teoría.

No obstante, los intentos de contrastar empíricamente estas teorías no han resultado muy positivos.\footnote{%
    Como resumen COURANT et al. (1986, pp. 279–280): \textit{(...) Pese a toda su elegancia y racionalidad, el modelo del ciclo vital no ha funcionado muy bien en las pruebas empíricas... Tampoco los intentos de contrastarlo con microdatos de corte transversal han resultado especialmente satisfactorios.}
} Se han propuesto diversas modificaciones para adaptar la teoría a los datos: introducir un motivo de herencia (recuérdese BARRO y su modelo de horizonte infinito), suponer imperfecciones en los mercados de capitales (modelos con restricciones de liquidez), asumir que la función de utilidad del consumo cambia con el tiempo (inconsistencia temporal de preferencias) o especificar alguna forma concreta de expectativas sobre la renta futura (modelos de cuasi-racionalidad o racionalidad limitada). Estas modificaciones aparecen ad hoc, ya que se requieren supuestos distintos para explicar cada anomalía empírica. 

Presentamos ahora una de las alternativas más completas proveniente de la corriente conductista de la economía; así como también algunos hechos empíricos especialmente relavantes a la hora de analizar el consumo. Empecemos por lo teórico.

\subsection{Behavioral Life Cycle: SHREFFIN y THALER (1989)}
Uno de los trabajos más interesantes es el propuesto por SHREFFIN y THALER (1989): la Hipótesis del Ciclo Vital Conductual (Behavioral Life Cycle, BLC). 

Estos autores argumentan que, en los modelos anteriores, muchos factores relevantes quedan descartados. Por ejemplo, el momento en que se percibe la renta a lo largo de los años o dentro de un mismo año (siempre que existan mercados de capital eficientes), así como la forma que adopta la riqueza (capital humano frente al valor del patrimonio inmobiliario). Factores que, a su parecer, son en realidad determinantes a la hora de modelizar la demanda de consumo corriente.

Sin ánimo de desarrollar en extenso el modelo, presentaremos sus supuestos principales así como algunas conclusiones y predicciones que realizan con el objetivo de ser contrastadas.

El modelo propuesto se basa en un problema de optimización intertemporal del consumo, de horizonte finito y con las particularidades/características de un modelo propio de la Economía de la Información: el \textbf{modelo Principal-Agente}. 

\subsubsection{Características conductuales: ¿cuáles son los intereses contraspuestos?}
Los autores plantean la existencia de un agente \textbf{Planificador} y un agente \textbf{Ejecutor} con intereses contrapuestos: el planificador vive toda la vida por lo que su intención es alisar lo máximo posible el consumo, mientras que el ejecutor se renueva cada periodo y su utilidad se deriva de la maximización del consumo en éste.

Más aún, para definir el problema de cada agente, plantean tres \textbf{conductas o características conductuales} que se reflejarán en su relación.

\paragraph{Problema económico: autocontrol y tentación}
En primer lugar, el autocontrol y su contraparte, la tentación. 

El autocontrol es necesario porque el consumo inmediato es siempre una alternativa atractiva frente a posponer el consumo.

El autocontrol es costoso y los agentes económicos (humanos) utilizan diversos mecanismos (como planes de pensiones o reglas heurísticas) para hacer frente a la dificultad de distribuir interteporalmente su consumo (especialmente, posponer una parte significativa de su consumo hasta la jubilación). De hecho, esta especificación no dista mucho de las observaciones que arrojado la psicología. El destacado psicólogo W. JAMES (1981)  señala que el atributo clave de las decisiones de autocontrol es la \textit{sensación de esfuerzo} que está presente.\footnote{%
  \textit{El esfuerzo de la atención es, por tanto, el fenómeno esencial de la voluntad. Todo lector debe saber por su propia experiencia que esto es así, pues todo lector habrá sentido el dominio de alguna pasión intensa. ¿En qué consiste la dificultad para una persona dominada por una pasión imprudente de actuar como si esa pasión fuera sensata? Desde luego, no hay ninguna dificultad física. Es tan fácil físicamente evitar una pelea como iniciarla, guardar el dinero como gastarlo en caprichos, alejarse como acercarse a la puerta de una coqueta. La dificultad es mental: consiste en mantener ante la mente la idea de la acción sensata.}
} 

Por tanto, el término autocontrol implica que los intercambios entre gratificación inmediata y beneficios a largo plazo conllevan un conflicto que no está presente en una elección entre, por ejemplo, una camisa blanca y una azul. Es decir, el autocontrol conlleva unos costes de \textbf{esfuerzo} reseñables.

Para captar formalmente la idea de este conflicto, los autores emplean una estructura de preferencias dual:\footnote{%
    Consideremos un individuo cuya vida se extiende durante $T$ periodos, siendo el último el de la jubilación. La secuencia de renta vital es $y=(y_1,\dots,y_T)$. Para simplificar, suponemos un mercado de capitales perfecto y un tipo de interés real nulo. La renta en la jubilación $y_T$ es cero:
\begin{la}
    \item \textbf{Ejecutor}. Dado que la tentación depende de las oportunidades de consumo inmediato, definimos un conjunto de oportunidades $X_t$ que representa las elecciones factibles de consumo en la fecha $t$. Si fuera libre de elegir en este conjunto, el ejecutor miope seleccionaría el valor máximo factible de $c_t$ (ya que eso maximiza U_t sobre $X_t$);
    \item \textbf{Planificador}. El planificador, en general, preferiría un $c_t$ menor. Supongamos que el planificador desea reducir el consumo ejerciendo fuerza de voluntad. Asumimos que, si el ejercicio de la fuerza de voluntad reduce $c_t$, debe existir algún coste psíquico. De no ser así, ejercer autocontrol no implicaría esfuerzo y problemas como el exceso de gasto o de consumo no existirían.
\end{la}
}

\eqblock{
\begin{aligned}
    \text{Planificador}:\;\;\succsim_P\qquad \underset{\substack{\text{\scriptsize Coexistentes +}\\\text{\scriptsize mutuamente inconsistentes}}}{\perp}\qquad \text{Ejecutor}:\;\;\succsim_E\\
    \text{Características}:
    \begin{aligned}
        &\uparrow\text{Fuerza de voluntad}\iff\downarrow c_t\\
        &\uparrow\text{Fuerza de voluntad}\iff \downarrow u_t\\
        &\uparrow \uparrow\text{Fuerza de voluntad}\iff \downarrow\downarrow c_t\\
        &\{\searrow\text{Fuerza de voluntad}+\searrow c_t\}\iff \uparrow T
    \end{aligned}
\end{aligned}
}{Característica conductual (1): Autocontrol y esfuerzo (fuerza de voluntad, tentación). Hipótesis Conductual de los Ciclos Vitales}

Uno orientado al largo plazo y otro al corto plazo. El ejecutor tiene control directo sobre la elección de consumo ($c_t$) y, dado que el autocontrol es costoso, el planificador puede buscar diferentes mecanismos para alcanzarlo.

\paragraph{Reglas: Internas y externas}
La técnica propuesta por los autores para la solución al conflicto entre las preferencias del planificador y del ejecutor consiste en que el planificador restrinja las elecciones futuras imponiendo restricciones que alteren $X_t$ (conjunto de elección presente). Por ejemplo, depositar fondos en un plan de pensiones que no permite retiradas reduce la renta disponible y, por tanto, reduce el conjunto de elección del ejecutor.\footnote{%
    En un caso extremo (equivalente a la modelización expuesta en los apartados anteriores), si el planificador pudiera elegir una regla que comprometiera completamente el consumo futuro a una trayectoria determinada. Como los ejecutores no tendrían decisiones que tomar, no sería necesario ejercer fuerza de voluntad. En esta situación, el planificador elegiría el consumo que maximizará su Función de Utilidad Intertemporal $V$ sujeto a la restricción presupuestaria, manteniendo el nivel de esfuerzo en $0$. 
Por tanto, la hipótesis del ciclo vital puede interpretarse como un caso particular del modelo de ciclo vital conductual (BLC) en el que o bien los costes de autocontrol son nulos, o bien existe una regla de compromiso óptima. Las predicciones de ambos modelos divergen porque ninguna de estas condiciones es realista: las personas sin costes de autocontrol son excepcionales y las reglas de compromiso perfectas no suelen estar disponibles. Aunque existen planes de pensiones y otros instrumentos de ahorro, la oferta es limitada y no determina completamente un plan de consumo. Además, la incertidumbre sobre los ingresos y las necesidades de gasto hace que estos planes no sean plenamente operativos.
}

Analíticamente:

\eqblock{
\begin{aligned}
    X_t^*\subset X_t\iff \underset{\text{\scriptsize utiliza mecanismos con reglas}}{\text{Planificador}}:
    \left\{\begin{aligned}
        &\text{Externos}\\
        &\text{Internos}
    \end{aligned}\right\}\Longrightarrow\;\;\underset{\substack{\text{\scriptsize disminuyen esfuerzo requerido}\\\text{\scriptsize para mantener un nivel de consumo}}}{\downarrow\text{FV}(\bar c_t)}
\end{aligned}
}

No obstante, existen límites al tipo de reglas que pueden aplicarse con bajos costes de autocontrol.\footnote{%
    La literatura psicológica sobre control de impulsos (por ejemplo, Ainslie, 1975) sugiere que las reglas eficaces deben cumplir ciertas características: deben ser simples (ya que las respuestas complejas requieren atención consciente), deben tener excepciones bien definidas y poco frecuentes, y deben ser dinámicamente estables (los hábitos no cambian fácilmente). Tanto las reglas internas como las externas son, por tanto, soluciones de segundo orden; en consecuencia, los modelos descriptivos del ahorro deben reflejar estas soluciones de segundo mejor que adoptan los individuos reales.
}

\paragraph{Contabilidad mental: Framing effect}
Esto nos conduce a la última característica conductual, que no es más que una implicación directa de las dos anteriormente expuestas: la descomposición de la riqueza (humana y activos) en distintas cuentas, denominadas \textbf{cuentas mentales}. Esto origina lo que en terminología económica se conoce como \textbf{framing-effect}: el consumo (o ahorro) se ve directamente afectado por la cuenta en la que se encuadran o describen los incrementos de riqueza.\footnote{%
    Como ejemplo, los autores mencionan que el modelo predice que un incremento de renta en forma de pago único (bonus) será tratado de forma distinta a la renta regular, incluso si ese bonus es completamente anticipado.
}

\eqblock{
\begin{aligned}
    &\text{Contabilidad mental}\Longrightarrow\text{Framing-effect}\\[1em]
    &\underset{\downarrow C_t\Longrightarrow\;\downarrow\downarrow\downarrow\text{FA}}{\text{Renta corriente ($I$)}}\quad |\quad \underset{\downarrow C_t\Longrightarrow\downarrow\downarrow\text{FA}}{\text{Activos actuales ($A$)}}\quad |\quad \underset{\downarrow C_t\Longrightarrow\downarrow\text{FA}}{\text{Renta futura ($F$)}}
\end{aligned}
}{}

Una versión simple divide la riqueza en tres componentes: renta corriente (I), activos actuales (A) y renta futura (F). En el modelo BLC, la propensión marginal a consumir depende de la cuenta en la que se clasifique la riqueza. Esto contrasta con el modelo tradicional del ciclo vital, que trata la riqueza como completamente fungible en mercados de capitales perfectos. En la teoría estándar, la propensión marginal a consumir es la misma ante distintos incrementos de riqueza: un bonus de 1000 dólares, un premio de lotería de 1000 dólares, un aumento de 1000 dólares en el valor de la vivienda o una herencia futura con valor presente de 1000 dólares. En cambio, el modelo conductual supone que los hogares clasifican estos incrementos en distintas cuentas mentales, algunas más “tentadoras” que otras.







Aunque las reglas internas de los hogares son contingente y dependientes del contexto, existen suficientes elementos comunes como para generar predicciones agregadas útiles. Uno de los elementos más importantes es 

La teoría BLC postula desigualdades específicas en la propensión marginal a consumir según el origen de la riqueza, no arbitrarias, sino resultado de mecanismos adaptativos para facilitar el ahorro. Esta descomposición constituye un ejemplo de framing (Kahneman y Tversky, 1984). Mientras que la teoría tradicional asume invariancia al encuadre, el modelo BLC introduce dependencia del encuadre.

Para ilustrar el sistema de tres cuentas, consideremos un hogar que utiliza una regla de pensiones según la cual se deduce una fracción s de la renta en cada periodo, sin acceso a esos fondos hasta la jubilación. Los saldos serían:
	1.	Cuenta de renta corriente: I = (1 - s)y_t.
	2.	Cuenta de riqueza actual: $A = \sum_{i=1}^{t-1} \big[(1 - s)y_i - c_i\big]$.
	3.	Cuenta de riqueza futura: suma de ingresos futuros netos más la riqueza acumulada en el plan de pensiones.

Este esquema es una simplificación. En la práctica, los hogares subdividen aún más estas cuentas (por ejemplo, ahorro para educación). Además, existe ambigüedad en la clasificación de ciertos ingresos: los dividendos pueden tratarse como renta corriente, mientras que otras rentas del capital se asignan a activos. El patrimonio inmobiliario también puede clasificarse como riqueza futura o actual, según el hogar.

Aunque este sistema puede parecer extraño para un economista, es análogo al funcionamiento contable de muchas universidades privadas, que distinguen entre fondos corrientes y fondos patrimoniales (endowment), con reglas estrictas sobre su uso.

Si el individuo desea ahorrar más allá de lo permitido por el plan de pensiones, debe recurrir al ahorro discrecional, lo que requiere esfuerzo de autocontrol para generar ese ahorro adicional, no gastarlo antes de tiempo y evitar endeudarse contra ingresos futuros. El coste de ese esfuerzo depende inversamente de la tentación. En el modelo, la tentación de gastar un dólar adicional depende de su ubicación en las cuentas mentales: es mayor en la renta corriente, menor en los activos actuales y aún menor en la riqueza futura.



 La contraparte del autocontrol sería entonces la \textbf{tentación}, ya que algunas situaciones son menos propicias para el ahorro que otras















Apoyándonos en la investigación realizada sobre estos temas por psicólogos y otros científicos sociales (véanse, por ejemplo, Ainslie, 1975, y Mischel, 1981), somos capaces de formular predicciones específicas sobre cómo el comportamiento real de ahorro de los hogares difiere del modelo idealizado de ciclo vital.

Nuestro modelo produce tres proposiciones que están significativamente en conflicto con la hipótesis del ciclo vital en este ámbito general. En esta sección se analiza la sensibilidad del consumo a la renta. Las dos secciones siguientes se refieren a los casos especiales de bonificaciones y ganancias inesperadas (windfalls). 



\subsection{Estancamiento secular: KEYNES} 
Según la teoría de Keynes, ante un crecimiento de renta sostenido, la propensión media al consumo se reduce, y aumenta la tasa de ahorro. Por ello, para mantener el equilibrio de pleno empleo, es necesario que el porcentaje sobre PIB de los otros componentes de la demanda agregada aumente. No obstante, en tanto la inversión depende principalmente de los \textit{animal spirits}, lo preferible será aumentar el gasto público. Basándose en esto, algunos economistas predijeron que la economía experimentaría lo que llamaron \textbf{estancamiento secular} (Hansen, 1939), entendida como una larga depresión de duración indefinida, a menos que el sector público utilizara la política fiscal para incentivar la demanda agregada. 

SUMMERS (2014), resucitó estas conjeturas y las aplicó a la situación económica actual de los países avanzados (Japón, Estados Unidos y Europa), afirmando que: \textit{Puede ser imposible que la economía alcance simultáneamente el pleno empleo, un crecimiento satisfactorio y la  estabilidad financiera simplemente mediante la operación de la política monetaria convencional}. 

Summers justificaba su observación en que durante los años de bonanza desde el pinchazo de la burbuja Puntocom a principios de los años 2000's y el comienzo de la crisis subprime en 2007, la economía de Estados Unidos, aún teniendo unos tipos de interés muy bajos y generando una burbuja en los mercados de activos de enormes proporciones, no experimentó presiones inflacionistas, no logró bajar la tasa de desempleo sustancialmente por debajo de la media histórica y tampoco logró crecimientos significativamente por encima de la tendencia. 

Desde una perpectiva defensiva, EICHENGREEN (2015) aclaró y distinguió entre tres explicaciones coherentes dentro del marco ortodoxo para el estancamiento secular:
\begin{lnum}
    \item \textbf{Tendencias de ahorro mundial}. En primer lugar, resulta necesario considerar el impacto de los cambios demográficos sobre las decisiones de ahorro y consumo: dado que el ahorro de las familias varía a lo largo del ciclo vital, cambios en la estructura por edades de la población producen un efecto composición. Así, por ejemplo, cuando la generación del \textit{baby boom} estaba en los años medios de su vida laboral se registró un aumento del ahorro agregado y, por las mismas razones, en el futuro con la disminución del peso relativo de la población joven se produciría una disminución. Asimismo, el aumento de la longevidad debería llevar asociado un incremento del ahorro. Otro factor que impulsaría el ahorro sería el pesimismo sobre las expectativas de crecimiento futuro del PIB asociado a un crecimiento reducido de la productividad y a la disminución de la población en edad de trabajar. Como ya se ha señalado, menores rentas futuras llevan a un mayor ahorro precautorio para suavizar intertemporalmente el consumo;
    \item \textbf{Inversión: Población, tecnología y productividad}. En segundo lugar, por lo que respecta a la inversión, un menor crecimiento de la población y de la productividad hacen esperar una menor demanda futura que justifique la acumulación de bienes de inversión para la producción futura. Por otra parte, un menor crecimiento de la población en edad de trabajar también requiere de un menor crecimiento de la inversión para mantener una determinada relación capital-trabajo; y,
    \item \textbf{Precios relativos de inversión}. En tercer lugar, la disminución tendencial de los precios relativos de los bienes de inversión hace que para acumular un determinado nivel de capital sea necesario un volumen de inversión menor. Es decir, el menor precio de los bienes de inversión significa que un determinado nivel de ahorro puede comprar mucho más capital que antes. Por ejemplo, GORDON (1990) hace hincapié en el papel de la investigación y el desarrollo que se materializa en bienes de inversión más baratos y eficientes.
\end{lnum}

En definitiva, los cambios demográficos y tecnológicos sugieren que puede producirse un déficit de demanda de consumo y de inversión junto con una cierta presión hacia un aumento del gasto público. No obstante, si bien esta es una caracterización global de la situación de la mayoría de los países avanzados, no todos ellos sufrirán los cambios demográficos ni la desaceleración de la productividad con la misma intensidad, y la situación de partida por lo que respecta a las finanzas públicas también son diferentes internacionalmente. 

El contexto económico asociado a un régimen de estancamiento secular tiene consecuencias importantes para las políticas económicas. En relación a la política monetaria, su eficacia para estimular el consumo y la inversión es muy reducida cuando los tipos de interés iniciales están cerca del ELB. Este déficit de demanda nacional se manifestaría en movimientos internacionales de capital que podrían acentuar los desequilibrios exteriores. En principio, cabría esperar que los países abocados a un mayor envejecimiento exportarán capital, de manera que aumentarían los superávits de sus balanzas por cuenta corriente, mientras que en los países con poblaciones más jóvenes tendrían balanzas por cuenta corriente deficitarias. Esta salida de capitales hacia los mercados emergentes llevaría asociada una disminución de los tipos de cambio reales para los países industriales, una mayor competitividad y una mayor demanda de exportaciones. Adicionalmente, cabe esperar migraciones internacionales desde los países con poblaciones más jóvenes hacia aquellos otros con poblaciones más envejecidas, si bien todos los países experimentarán una escasez relativa de trabajadores jóvenes.

Para hacer frente a estos impactos negativos del estancamiento secular, SUMMERS plantea dos estrategias:\footnote{%
    En su artículo \textit{U.S. Economic Prospects: Secular Stagnation, Hysteresis, and the Zero Lower Bound} SUMMERS considera una tercera estrategia basada en la paciencia: en ocasiones las políticas económicas tienen limitado impacto y esto puede revertirse en el futuro. El autor sugiere que esta parece ser la estrategia seguida por países como Japón.
} 
\begin{lnum}
    \item Buscar alternativas para reducir aún más las tasas de interés reales. Por ejemplo, estableciendo un objetivo de tasa de inflación más alto. No obstante, estas estrategias conllevan el riesgo de que incluso si aumentan el nivel de demanda, también es probable que aumenten los riesgos para la estabilidad financiera, lo que a su vez puede tener consecuencias negativas sobre la demanda; y, 
    \item Aumentar la demanda mediante el incremento de la inversión y la reducción del ahorro. Por ejemplo, aumentando la inversión pública, reduciendo las barreras sobre la inversión privada o fomentando las exportaciones.
\end{lnum}
 










\clearpage
\section{Conclusión} 

\subsection{Extensiones y relación con otras partes del Temario}


\subsection{Opinión}



\subsubsection{Idea final (Salida o cierra)}


\clearpage
\pagenumbering{gobble}
\MyBibliography
\MyBibItem{Autor}{Título}{Año}{DOI -- LINK}



\clearpage
\anexo{Aproximación ortodoxa (dinámica): Derivación analítica de los modelos}\label{anx:derivacion_modelos}

\textcolor{magenta}{Breve introducción}

\textbf{1.- Hipótesis de la Renta Permanente (FRIEDMAN, 196X)}


La hipótesis de la renta permanente hacía hincapié en las variaciones \textbf{estocásticas} de la renta a lo largo del tiempo. 

Así, el modelo propuesto por FRIEDMAN descompone al consumo en dos partes: permanente ($C_t^P$) y temporario ($C_t^T$); como consecuencia de existir dos tipos de renta: permanente ($Y_t^P$) y temporario ($Y_t^T$). ¿Cómo se relacionan entre ellos? 

Para FRIEDMAN, la renta permanente es una porción de la riqueza total del individuo: $Y_t^P=r\cdot W_t$. Puede entenderse como el flujo de renta que puede ser consumido en un periodo de tiempo \textbf{sin afectar al valor de la fuente de la cual proviene}. Por ello, se concibe (en su formulación más básica), como el rendimiento obtenido en el periodo $t$ por la riqueza acumulada por el individuo hasta ese momento ($W_t$), compuesto por activos reales, activos financieros (dinero, bonos), y riqueza humana siendo. No es de extrañar, por tanto, que la estimación empírica de la riqueza se convirtiera en una cuestión central en la especificación de la función de consumo de Friedman; así como también en su principal debilidad. 

Por otro lado, FRIEDMAN concibe la renta transitoria como una variable compuesta por dos elementos: (i) shocks transitorios que afectan a la renta; y, (ii) incrementos contingentes de renta que serán integrados como permanentes. En ambos casos, se formulan como shocks estocásticos.

Respecto al consumo, FRIEDMAN postula la existencia de un consumo permanente proporcional a la renta permanente. Este coeficiente proporcional, por ejemplo $k$, depende de factores tales como el tipo de interés, los gustos y elementos demográficos. Y, por otro lado, el gasto de un periodo cualquiera ($C_t$), se descompone en un elemtno permanente ($C_t^P$), que ya hemos definido, y un componente transitorio. ¿Por qué? Por ejemplo, este elemento transitorio puede racionalizarse como periodo de compras extraordinarias, gastos de enfermedad, etc. 

Con esta definición, podemos ver como FRIEDMAN enfatiza la noción de redistribución intertemporal del consumo: el agente \textit{siempre quiere} consumir prácticamente la misma porción, y son los sucesos contingentes lo que le separan de esta tendecia lineal. 

Hemos definido el comportamiento individual de una persona cualquier. Pero, lo que FRIEDMAN pretende es dar respuesta a los patrones de consumo agregado, por tanto: ¿cómo podemos agregar estos patrones individuales de consumo?

Para conseguir la agregación debemos añadir una serie de supuestos. En términos individuales tenemos:

$\left.\begin{aligned}
            \begin{aligned}
                &\text{Renta}:Y_{i,t}^P=r\cdot W_{i,t}\quad \text{y,}\quad Y_{i,t}=Y_{i,t}^P+Y_{i,t}^T\\[1em]
                &\text{Consumo}:C_{i,t}^P=k\cdot Y^P_{i,t}\quad \text{y,}\quad C_{i,t}=C_{i,t}^P+C_{i,t}^T
            \end{aligned}\quad ;\quad \forall i\in I
  \end{aligned}\hspace{1em}\right.$


Por lo tanto, para agregarlos supondremos que los shocks transitorios de consumo se cancelan entre sí: los individuos tendrán shocks positivos y otros negativos, aunque a nivel individual puedan ser muy importantes. \textcolor{magenta}{¿Por qué?} (intuición económica de FRIEDMAN, por qué es importante esto para FRIEDMAN)

En segundo lugar, supondremos que, a nivel agregado, las rentas transitorias positivas y negativas de los agentes no se cancelan en su totalidad. Es decir, existe un sesgo sistemático que empuja la renta de la economía hacia arriba:

$\sum_i Y_{i,t}\Longrightarrow Cov(Y_t^P,Y_t^T)=0\quad \text{y},\quad \underset{$\exists$\text{\scriptsize  sesgo sistemático en la renta}}{\mathbb E[Y_t^T]\neq 0}$


En tercer, y último lugar, asumiremos que los componentes del consumo, transitorios y permanentes, no están correlacionados entre si.

$\sum_iC_{i,t}\Longrightarrow Cov(C_t^P,C^T_t)=0\quad \text{y,}\quad \underset{\text{\scriptsize $\not\exists$ sesgo sistemático en el consumo}}{\mathbb E[C_t^T]=0}$

Evidentemente, la intención de FRIEDMAN al formular estas dos últimas hipótesis es reconciliar la Hipótesis de la Renta Permanente con las observaciones empíricas. 

La idea central es que los shocks transitorios de renta son accidentales o aleatorios, pero no necesariamente simétricos entre individuos o periodos, por lo que al agregarlos no tienen por qué compensarse exactamente y de este modo SIEMPRE habrá un componente transitorio en la renta observada (sensu contrario, la renta observada NUNCA será igual a la renta permanente).

La lógica contraria respecto al consumo es precisamente lo que da lugar al fenómeno de \textbf{suavización del consumo}. ¿Por qué? Si a nivel agregado el componente del gasto de cancela entre sí (probablemente), entonces el consumo observado depende fundamentalmente de la renta permanente y, en consecuencia, los cambios transitorios de renta no generan sistemáticamente cambios en el consumo. Más bien, los agentes utilizarán estos cambios transitorios de renta precisamente para suavizar su consumo a lo largo del tiempo.

Analíticamente:

$Cov(C_t^T,Y_t^T)=0$

Con estos supuestos estadísticos adicionales para reconciliar el modelo con la evidencia empírica,  FRIEDMAN es capaz de reconciliar su teoría con los datos de sección cruzada (KUZNETS y DUESENBERRY, [\authorfont{Ver Anexo \ref{anx:antecedentes_historicos}}]), que muestran unas relaciones de consumo-renta más altas en los tramos de renta más bajos. 


\textbf{Teoría del Ciclo Vital (ANDO y MODIGLIANI, 196X)}

La Hipótesis del Ciclo Vital se centra en las variaciones \textbf{predecibles} de la renta a lo largo del ciclo vital para explicar las fluctuaciones en el consumo.

Para elaborar su teoría, utilizan para el modelo de FISHER sobre decisiones intertemporales con el fin de determinar el consumo de una persona en función de la renta/riqueza que prevee percibir a lo largo de toda su vida. 

En primer lugar, se define el consumo de un individuo cualquier como:

$\begin{aligned}
  C_{i,t}=K_{i,t}^T\cdot W_{i,t},\qquad \text{donde,}
    \left\{\begin{aligned}
        &C_{i,t}\equiv \text{Consumo en $t$ de $i$}\\[0.5em]
        &W_{i,t}^T\equiv \text{Riqueza, descontada a $t$, de $i$}\\[0.5em]
        &\underset{\text{\scriptsize $Cov(K_{i,t}^T,y_{i,t})=0\,$ no depende de la renta}}{K_{i,t}^T\equiv \text{Proporción que depende de diversos factores}}
    \end{aligned}\right.
\end{aligned}$

\textcolor{magenta}{Es decir, el consumo se define como función de ... (continuar)}

La riqueza del individuo en el momento actual, por otro lado, se define como:

$\begin{aligned}
  W_{i,t}=\underbrace{a_{i,t-1}}_{\text{\scriptsize riqueza no humana}}+\underbrace{\sum_{j=t}^T\frac{\mathbb E_t[y_{i,{j}}]}{(1+r)^{j-t}}}_{\text{\scriptsize riqueza humana $:=\,H_{i,t}$}}
\end{aligned}$

\textcolor{magenta}{Es decir, la riqueza se define como como función de...}

Juntando ambas ecuaciones, podemos definir el consumo como:

$C_{i,t}=K_{i,t}^T\cdot(a_{i,t-1}+h_{i,t})$

\textcolor{magenta}{Es decir, como ... (continuar)}

\textcolor{magenta}{Por último, para hacer esta ecuación operativa en el contraste empírico debemos hacer una serie de cambios .... (continuar, ¿por qué?)} Con todo, la ecuación resultante queda:

$\begin{aligned}
  \overset{\substack{\text{\scriptsize para}\\\text{\scriptsize hacerlo}\\\text{\scriptsize operativo}}}{\Longrightarrow}\quad 
    \text{si,}\,\underset{\text{\scriptsize o, más generalmente, $\mathbb E_t[y_{i,t+s}]$ sigue HEA}}{\mathbb E_t[y_{i,t+s}] = \phi_s\,y_{i,t},\,\,\forall s}\Longrightarrow W_{i,t}=a_{i,t-1}+\Phi_{i,t}y_{i,t}
\end{aligned}$

Por tanto, el consumo queda:

$C_{i,t}=K_{i,t}^T(a_{i,t-1}+\Phi_{i,t}y_{i,t})$

Tenemos ya el consumo del agente en función de su riqueza y determinación de la riqueza vital.


Igual que antes, debemos agregar los patrones de consumo individual de los agentes, de acuerdo con esta Teoría, con el fin de obtener el consumo agregado. 

Imponiendo una serie de parámetros adicionales respecto a las distribuciones de determinadas variables (propensiones medias al consumo, distribuciones de riqueza y renta y estructura demográfica), podremos obtener la siguiente función agregada de consumo.

\textcolor{magenta}{En primer lugar, continuar aquí.... (explicar qué significa y porqué es importante)} . Por tanto, esta condición se cumple si imponemos la siguiente condición:

$\begin{aligned}
  \mathbb E[K_{i,t}^T\;|\;T_i=\tau]=K_{t}^T\quad &&\underset{\text{\scriptsize grupo de edad}}{\forall\tau\in[0,T]}
\end{aligned}$

\textcolor{magenta}{En segundo lugar, continuar aquí.... (explicar qué significa y porqué es importante)} . Por tanto, esta condición se cumple si imponemos la siguiente condición:

$\begin{aligned}
  \underset{\text{\scriptsize distribución de renta para cada grupo de edad}}{F_{Y|T}(y\;|\;\tau)=\mathrm P[y_{i,t}\leq y\;|\;T_i=\tau]}\;\;\wedge\;\; \underset{\text{\scriptsize distribución de riqueza para cada grupo de edad}}{F_{A|T}(a\;|\;\tau)=\mathrm P[a_{i,t}\leq A\;|\;T_i=\tau]}\quad &&\perp\forall \;\;T
\end{aligned}$


\textcolor{magenta}{En tercer lugar, continuar aquí.... (explicar qué significa y porqué es importante)} . Por tanto, esta condición se cumple si imponemos la siguiente condición:

$\begin{aligned}
  \underset{\text{\scriptsize pirámide poblacional}}{F_T(t)=\mathrm P[T_i\leq \tau]}&&\perp \forall t
\end{aligned}$

Siempre y cuando se cumplan estas condiciones, podremos obtener una función de consumo con la siguiente forma funcional:

$\begin{aligned}
  C_t\overset{(\checkmark)}{=}\sum C_{i,t}=\sum_iK_{i,t}^Ta_{i,t-1}+\sum_{i}K_{i,t}^T\Phi_{i,t}y_{i,t}\qquad \text{definimos,}
    \left\{\begin{aligned}
        &A_{t-1}&&\equiv\sum_ia_{i,t-1}\\
        &Y_t&&\equiv \sum_iy_{i,t}
    \end{aligned}\right.
\end{aligned}$

\textcolor{magenta}{Es decir, ... (continuar aquí explicando la ecuación y por qué las definiciones)}

La ecuación anterior es equivalente a:

$\begin{aligned}
  C_t=\underbrace{\frac{\sum_iK_{i,t}^Ta_{i,t-1}}{\sum_ia_{i,t-1}}}_{\text{\scriptsize proporción de consumo sobre riqueza total, $\alpha$}}A_{t-1}+\underbrace{\frac{\sum_iK_{i,t}^T\Phi_{i,t}y_{i,t}}{\sum_iy_{i,t}}}_{\text{\scriptsize proporción de consumo sobre renta total, $\beta$}}Y_{t}
\end{aligned}$

\textcolor{magenta}{Es decir, .... (continuar aquí con la descripción)}. En consecuencia, podemos reformular la ecuación como:

$C_t=\alpha \,A_{t-1}+\beta\, Y_t$

Obteniendo así la función de consumo agregada de la economia como \textcolor{magenta}{función de.... (continuar)}


\textbf{Modelo de HALL: Paseo aleatorio}

HALL parte de un problema de optimización intertemporal clásico a partir del cual podemos derivar la ecuación de EULER relativa al consumo: la condición de primer orden que resuelve el problema dinámico. Analíticamente, empezamos planteando el problema del consumidor:

$\begin{aligned}
      &\max U(0)=\sum_{t=0}^T u(c_t)\beta^t\\[0.5em]
      &\text{s.a.}\quad 
      \left\{\begin{aligned}
            &A_t=(1+r)A_{t-1}+y_t-C_t\\
            &A_0=\bar A>0\\
            &A_T=0
      \end{aligned}\right.
\end{aligned}$

Resolviendo, obtenemos:

$\begin{aligned}
  \overset{(\beta_t=\frac{1}{1+\rho})}{\Longrightarrow} \text{CPO}:\frac{\mathrm d u(C_t)}{(1+r)}=\frac{\mathrm d u(C_{t+1})}{(1+\rho)}
\end{aligned}$

Es decir, el equilibrio dinámico se alcanza con la igualdad entre utilidades marginales de dos periodos correlativos, el corriente y el siguiente esperado. 

De esta forma, \textcolor{magenta}{si imponemos HER y consideramos... (desarrollar textualmente aquí y analíticamente en abajo en la fórmula, con pasos, no texto sobre flecha, quita eso)} . Obtendríamos:

$\begin{aligned}
  \overset{\substack{\text{\scriptsize Aplicando HER +}\\\text{\scriptsize $r=\rho$}}}{\Longrightarrow}\quad \boxed{u'(C_t)=\mathbb E_t[u'(C_{t+1})]}
\end{aligned}$

Resuelto el problema de optimización, obtenemos el consumo óptimo en el marco biperiodo.

Sin embargo, nótese como esto no nos proporciona información adicional a la hora de llevar a cabo predicciones y/o contrastaciones empíricas ya que ello requeriría establecer una relación biunívoca entre \textbf{consumos y no entre utilidades derivadas del consumo}. 

Para alcanzar este punto, HALL asume (en la versión canónica) que la Función de Utilidad instantanea tiene una forma funcional cuadrática:

$\begin{aligned}
  u(C_t)=aC_t-\frac{b}{2}C_t^2
\end{aligned}$

Por tanto, si \textcolor{magenta}{derivamos...}, obtendremos:

$u'(C_t)=a-bC_t$

Sustituyendo la ecuación anterior en la CPO, obtendríamos:

$\begin{aligned}
  u'(C_t)=\mathbb E_t[u'(C_{t+1})]
\end{aligned}$

Lo que implica que:

$a-bC_t=\mathbb E_t[a-bC_{t+1}]$

Por las propiedades de laesperanza matemática, obtenemos que:

$C_t=\mathbb E_t[C_{t+1}]$

Es decir, que el mejor predictor del consumo futuro, es el consumo presente. Esta es la conocida expresión del modelo de HALL: el \textbf{paseo aleatorio}. Como vemos, el consumo presente es igual al valor esperado del consumo futuro.

Por tanto, para obtener una ecuación con la que podamos trabajar con datos lo único que debemos hacer es invertir la condición definiendo el consumo futuro no anticipado (HER): el consumidor elige $C_t$ en funcion de la riqueza total esperada, es decir, activos mas valor presente de rentas futuras esperadas. Si llega nueva informacion en $t+1$ que revisa al alza o a la baja la renta futura esperada, cambia la riqueza permanente y por tanto cambia el nivel optimo de consumo. 

Analíticamente:

$C_t=\mathbb E_t[C_{t+1}]\Longrightarrow C_{t+1}=C_t+\eta_{t+1}$

\textcolor{magetna}{Teniendo en cuenta que... ()qué es, hacer breve descripción de lo qué significa}

$\underset{\substack{\text{\scriptsize El consumo es una función de la renta,}\\\text{\scriptsize preservan la idea de FRIEDMAN}}}{C:Y\in \mathbb R^+\longrightarrow \mathbb R^+}$

Por tanto, obtenemos que:

$\begin{aligned}
  \text{Resultados HALL:}
  \left\{\begin{aligned}
        &\text{Expresado en $C$}:&&C_{t+1}=C_t+\eta_{t+1},&& \eta\sim N(0,\sigma_\eta)\\[0.5em]
        &\text{Expresado en $Y$}:&&C_{t+1}=\beta_0(y_{Dt}+y_{Lt})+\beta_1\,y_{st},&& y_{Lt}\sim N(0,\sigma_L)\quad \text{y,}\quad y_{st}\sim AR(1)
  \end{aligned}\right.
\end{aligned}$

Obteneiendo así las regresiones necesarias para llevar a cabo la contrastación empírica.

Por tanto, el modelo de HALL permite reconciliar la Hipótesis de la Renta Permanente y la Hipótesis del Ciclo Vital, unificándola bajo un mismo paraguas e incorporando la HER. Esto se deriva de la aportación fundamental de HALL y MISHKIN, separando la renta de los hogares en tres componentes: (i) un \textbf{componente determinístico} $y_{Dt}$, que aumenta con la edad hasta justo antes de la jubilación; (ii) un \textbf{componente estocástico} $y_{Lt}$, que fluctúa a medida que cambian las perspectivas de ingresos a lo largo de la vida y que se especifica como un paseo aleatorio (random walk); y, (iii) un \textbf{componente estocástico estacionario} $y_{st}$, que fluctúa debido a influencias transitorias y se describe mediante un proceso de media móvil. 

En síntesis, el modelo de HALL concluye que: (i) para predecir el consumo futuro ($t+1$) solo necesitamos el consumo presente ($t$); (ii) los valores retardados de cualquier otra variable que en principio pudiesen ser relevantes para predecir el consumo futuro ($t+1$) no tienen valor explicativo adicional una vez que consideramos el consumo presente ($t$); (iii) el término de error refleja la nueva información sobre rentas futuras, información que no estaba disponible en el presente ($t$).

\textbf{Consumo de bienes duraderos}

La diferencia clave entre los bienes no duraderos y los duraderos es que los primeros son una variable flujo, mientras que los segundos son una variable stock que proporciona un flujo de servicios (utilidad) en cada periodo de vigencia. De esta forma, deben ser incluidos de manera explícita en el problema de maximización del consumidor. 


Analíticamente, empezamos formulando el problema de maximización del consumidor:

$\begin{aligned}
        &\max U(0)=\sum_{s=0}^Tu(C_{t+s},D_{t+s})\\[0.5em]
        &\text{s.a.}\quad 
        \left\{\begin{aligned}
            &A_t=(1+r)A_{t-1}+y_t-C_t-\underset{\text{\scriptsize Inversión neta en stock de bienes durables}}{p_t^D(D_{t+1}-(1-\delta)D_{t})}\\
            &\bar A=a\\
            &A_T=0
        \end{aligned}\right.
\end{aligned}$

Para resolver el problema anterior debemos asumir dos condiciones:

$\begin{aligned}
  &u(C_t,D_t)=C_t^\alpha D_t^{1-\alpha}\\[0.5em]
  &\not\exists\text{Restricciones de crédito}
\end{aligned}$

\textcolor{magenta}{Asumiendo que ... (continuar, decir por qué y qué son)}. Con todo, podemos formular la relación óptima de composición del gasto en el Estado Estacionario:

$\begin{aligned}
  \frac{C_t}{p^D\,D_t}=\frac{\alpha}{1-\alpha}(r_{t+1}+\delta)
\end{aligned}$

\textbf{IMPORTANTE.-} Esta última condición debe leerse como una condicion intratemporal de composicion entre no duraderos y servicios del stock durable. No es una CPO de consumo agregado en el sentido tradicional, sino una igualdad entre utilidad marginal relativa y coste de uso del durable: $\frac{u_D}{u_C}=\frac{1-\alpha}{\alpha}\frac{C_t}{D_t}$

Esta nueva formulación tiene tres implicaciones principales.

En primer lugar, la \textbf{desvinculación consumo-gasto}. El consumo relevante no coincide con el gasto cuando hay bienes duraderos. Es decir, el consumo en no duraderos se eligen como flujo periodo a periodo, lo que corresode con el gasto. En cambio, los duraderos se eligen como stock y por tanto implican una decision de cartera intertemporal, desvinculando el \textit{consumo} del \textit{gasto}. Desde la óptica del investigador: el gasto en duraderos será mucho más volátil y más sensible a tipos de interés, renta esperada y precios relativos que el gasto en no duraderos. No porque el consumidor suavice peor, sino porque el \textbf{gasto en duraderos contiene un componente de inversión}

En segundo lugar, la relación de \textbf{composición óptima de consumo}. La composición óptima entre no duraderos y servicios del stock durable depende del coste de uso del durable ($p_t^D$), no solo de la renta. No obstante, los cambios sobre la renta permanente también tienen proyección sobre el modelo. En no duraderos, el efecto es claro: mayor riqueza permanente implica mayor flujo de consumo desde hoy. En duraderos, el hogar deseará un stock objetivo mayor. Para acercarse a ese nuevo stock objetivo, aumenta la compra de duraderos hoy; pero como el stock se acumula y se deprecia lentamente, el consumo será gradual incluso si el gasto pega un salto (e $t$). Por eso el gasto en duraderos puede sobrerreaccionar al principio (hipersensibilidad).

En tercer lugar, el \textbf{ajuste empírico}. Los shocks de renta o de tipos de interés pueden generar respuestas muy distintas en no duraderos y en gasto en duraderos, porque este ultimo es ajuste de stock. Tal y como se ha anticipado en las dos implicaciones anteriores. En particular, el mejor ajuste empírico se origina en: (i) las \textbf{perturbaciones en renta}, pues parte de la hipersensibilidad se origina en las observaciones de compras (gasto) de bienes duraderos, que no son un flujo puro de consumo sino un ajuste de stock; y, (ii) los \textbf{retardos}, pues incorporando duraderos aparecen de forma natural retardos de ajuste. Por tanto, no hace falta suponer miopía para explicar por que el gasto no se mueve uno a uno con la nueva riqueza permanente. Basta con reconocer que el stock durable heredado del pasado condiciona la elección presente.

De esta forma, trabajar con no duraderos y servicios permitiría evitar que la dinámica de los bienes duraderos contamine la contrastación del modelo básico. Una renta transitoria puede provocar una reacción fuerte del gasto en duraderos aunque el modelo intertemporal sea correcto. 

Se podría decir que la incorporación de bienes duraderos interpreta que si el parámetro de contrastación de HALL ($\beta_1$) es elevado, demuestra que la variable observada contiene un componente que no sigue la formulación básica para no duraderos. Lo que no contradeciría el fenómeno de la suavizacion del consumo; sino que explicaría por qué el gasto observado es \textit{hipersensible}, mezclando consumo y acumulación de stock durable. 

De hecho, los hogares pueden reducir temporalmente las compras de bienes duraderos cuando sus ingresos disminuyen, sin que se produzca una gran reducción de la utilidad a corto plazo. Cuando sus ingresos se recuperan, pueden reanudar sus compras para compensar las pérdidas anteriores en su stock de bienes duraderos.




\clearpage
\anexo{Aproximación conductista (dinámica): Derivación del modelo}\label{anx:conductista}

El modelo propuesto se basa en un problema de optimización intertemporal del consumo, de horizonte finito y con las particularidades/características de un modelo propio de la Economía de la Información: el \textbf{modelo Principal-Agente}. 

Los autores plantean la existencia de un agente \textbf{Planificador} y un agente \textbf{Ejecutor} con intereses contrapuestos: el planificador vive toda la vida por lo que su intención es alisar lo máximo posible el consumo, mientras que el ejecutor se renueva cada periodo y su utilidad se deriva de la maximización del consumo en éste.

Más aún, para definir el problema de cada agente, plantean tres \textbf{conductas o características conductuales} que se reflejarán en su relación.

\textbf{1.- Problema económico: autocontrol y tentación}

En primer lugar, el autocontrol y su contraparte, la tentación. 

El autocontrol es necesario porque el consumo inmediato es siempre una alternativa atractiva frente a posponer el consumo.

El autocontrol es costoso y los agentes económicos (humanos) utilizan diversos mecanismos (como planes de pensiones o reglas heurísticas) para hacer frente a la dificultad de distribuir interteporalmente su consumo (especialmente, posponer una parte significativa de su consumo hasta la jubilación). De hecho, esta especificación no dista mucho de las observaciones que arrojado la psicología. El destacado psicólogo W. JAMES (1981)  señala que el atributo clave de las decisiones de autocontrol es la \textit{sensación de esfuerzo} que está presente.\footnote{%
  \textit{El esfuerzo de la atención es, por tanto, el fenómeno esencial de la voluntad. Todo lector debe saber por su propia experiencia que esto es así, pues todo lector habrá sentido el dominio de alguna pasión intensa. ¿En qué consiste la dificultad para una persona dominada por una pasión imprudente de actuar como si esa pasión fuera sensata? Desde luego, no hay ninguna dificultad física. Es tan fácil físicamente evitar una pelea como iniciarla, guardar el dinero como gastarlo en caprichos, alejarse como acercarse a la puerta de una coqueta. La dificultad es mental: consiste en mantener ante la mente la idea de la acción sensata.}
} 

Por tanto, el término autocontrol implica que los intercambios entre gratificación inmediata y beneficios a largo plazo conllevan un conflicto que no está presente en una elección entre, por ejemplo, una camisa blanca y una azul. Es decir, el autocontrol conlleva unos costes de \textbf{esfuerzo} reseñables.

Para captar formalmente la idea de este conflicto, los autores emplean una estructura de preferencias dual.

La conducta del \textbf{ejecutor} se define entorno a la tentación. Dado que la tentación depende de las oportunidades de consumo inmediato, definimos un conjunto de oportunidades $X_t$ que representa las elecciones factibles de consumo en la fecha $t$. Si fuera libre de elegir en este conjunto, el ejecutor miope seleccionaría el valor máximo factible de $c_t$ (ya que eso maximiza $U_t$ sobre $X_t$);

Por el contrario, la conducta del \textbf{planificador} de define entorno al autocontrol. El planificador, en general, preferiría un $c_t$ menor. Supongamos que el planificador desea reducir el consumo ejerciendo fuerza de voluntad. Asumimos que, si el ejercicio de la fuerza de voluntad reduce $c_t$, debe existir algún coste psíquico. De no ser así, ejercer autocontrol no implicaría esfuerzo y problemas como el exceso de gasto o de consumo no existirían.

Con esto, consideremos un individuo cuya vida se extiende durante $T$ periodos, siendo el último el de la jubilación. La secuencia de renta vital es $y=(y_1,\dots,y_T)$. Para simplificar, suponemos un mercado de capitales perfecto y un tipo de interés real nulo. La renta en la jubilación $y_T$ es cero.

\textcolor{magenta}{Lo que está en eqblock hay que integrarlo en el texto (de forma que el texto de 1.- quede armónico) y definir formalmente las ecuaciones (si cabe). Tenerlo en cuenta}

\eqblock{
\begin{aligned}
    \text{Planificador}:\;\;\succsim_P\qquad \underset{\substack{\text{\scriptsize Coexistentes +}\\\text{\scriptsize mutuamente inconsistentes}}}{\perp}\qquad \text{Ejecutor}:\;\;\succsim_E\\
    \text{Características}:
    \begin{aligned}
        &\uparrow\text{Fuerza de voluntad}\iff\downarrow c_t\\
        &\uparrow\text{Fuerza de voluntad}\iff \downarrow u_t\\
        &\uparrow \uparrow\text{Fuerza de voluntad}\iff \downarrow\downarrow c_t\\
        &\{\searrow\text{Fuerza de voluntad}+\searrow c_t\}\iff \uparrow T
    \end{aligned}
\end{aligned}
}{Característica conductual (1): Autocontrol y esfuerzo (fuerza de voluntad, tentación). Hipótesis Conductual de los Ciclos Vitales}

Uno orientado al largo plazo y otro al corto plazo. El ejecutor tiene control directo sobre la elección de consumo ($c_t$) y, dado que el autocontrol es costoso, el planificador puede buscar diferentes mecanismos para alcanzarlo.

\textbf{2.- Reglas: Internas y externas}

La técnica propuesta por los autores para la solución al conflicto entre las preferencias del planificador y del ejecutor consiste en que el planificador restrinja las elecciones futuras imponiendo restricciones que alteren $X_t$ (conjunto de elección presente). Por ejemplo, depositar fondos en un plan de pensiones que no permite retiradas reduce la renta disponible y, por tanto, reduce el conjunto de elección del ejecutor.

Es decir, el planificador dispone de instrumentos \textbf{externos} (esrtategias) o \textbf{internos} (esfuerzo) para reducir el conjunto de consumo factible en el presente, haciendo que el óptimo sea menor que el óptimo sin restricción:

$X_t^*\subset X_t$

A su vez, la alternancia entre estos instrumentos hace que el esfuerzo requerido (desutilidad) para mantener este nivel de consumo restringido sea menor, por tanto optará por aprovechar estas estrategias:

$\underset{\substack{\text{\scriptsize disminuyen esfuerzo requerido}\\\text{\scriptsize para mantener un nivel de consumo}}}{\downarrow\text{FV}(\bar c_t)}$

Nótese que, en realidad, el caso extremo de esta formalización sería equivalente a la modelización expuesta en los apartados anteriores: el planificador puede elegir una regla que comprometiera completamente el consumo futuro a una trayectoria determinada. En este caso, como los ejecutores no tendrían decisiones que tomar, no sería necesario ejercer fuerza de voluntad. En esta situación, el planificador elegiría el consumo que maximizará su Función de Utilidad Intertemporal $V$ sujeto a la restricción presupuestaria, manteniendo el nivel de esfuerzo en $0$. 

Por tanto, la hipótesis del ciclo vital puede interpretarse como un caso particular del modelo de ciclo vital conductual (BLC) en el que o bien los costes de autocontrol son nulos, o bien existe una regla de compromiso óptima. Las predicciones de ambos modelos divergen porque ninguna de estas condiciones es realista: las personas sin costes de autocontrol son excepcionales y las reglas de compromiso perfectas no suelen estar disponibles. Aunque existen planes de pensiones y otros instrumentos de ahorro, la oferta es limitada y no determina completamente un plan de consumo. Además, la incertidumbre sobre los ingresos y las necesidades de gasto hace que estos planes no sean plenamente operativos.

No obstante, existen límites al tipo de reglas que pueden aplicarse con bajos costes de autocontrol. La literatura psicológica sobre control de impulsos (por ejemplo, AINSLIE, 1975) sugiere que las reglas eficaces deben cumplir ciertas características:
\begin{lnum}
  \item Deben ser simples (ya que las respuestas complejas requieren atención consciente), deben tener excepciones bien definidas y poco frecuentes, y deben ser dinámicamente estables (los hábitos no cambian fácilmente). Tanto las reglas internas como las externas son, por tanto, soluciones de segundo orden; en consecuencia, los modelos descriptivos del ahorro deben reflejar estas soluciones de segundo mejor que adoptan los individuos reales.
\end{lnum} 

\paragraph{Contabilidad mental: Framing effect}
Esto nos conduce a la última característica conductual, que no es más que una implicación directa de las dos anteriormente expuestas: la descomposición de la riqueza (humana y activos) en distintas cuentas, denominadas \textbf{cuentas mentales}. Esto origina lo que en terminología económica se conoce como \textbf{framing-effect}: el consumo (o ahorro) se ve directamente afectado por la cuenta en la que se encuadran o describen los incrementos de riqueza.\footnote{%
    Como ejemplo, los autores mencionan que el modelo predice que un incremento de renta en forma de pago único (bonus) será tratado de forma distinta a la renta regular, incluso si ese bonus es completamente anticipado.
}

\eqblock{
\begin{aligned}
    &\text{Contabilidad mental}\Longrightarrow\text{Framing-effect}\\[1em]
    &\underset{\downarrow C_t\Longrightarrow\;\downarrow\downarrow\downarrow\text{FA}}{\text{Renta corriente ($I$)}}\quad |\quad \underset{\downarrow C_t\Longrightarrow\downarrow\downarrow\text{FA}}{\text{Activos actuales ($A$)}}\quad |\quad \underset{\downarrow C_t\Longrightarrow\downarrow\text{FA}}{\text{Renta futura ($F$)}}
\end{aligned}
}{}

Una versión simple divide la riqueza en tres componentes: renta corriente (I), activos actuales (A) y renta futura (F). En el modelo BLC, la propensión marginal a consumir depende de la cuenta en la que se clasifique la riqueza. Esto contrasta con el modelo tradicional del ciclo vital, que trata la riqueza como completamente fungible en mercados de capitales perfectos. En la teoría estándar, la propensión marginal a consumir es la misma ante distintos incrementos de riqueza: un bonus de 1000 dólares, un premio de lotería de 1000 dólares, un aumento de 1000 dólares en el valor de la vivienda o una herencia futura con valor presente de 1000 dólares. En cambio, el modelo conductual supone que los hogares clasifican estos incrementos en distintas cuentas mentales, algunas más “tentadoras” que otras.







Aunque las reglas internas de los hogares son contingente y dependientes del contexto, existen suficientes elementos comunes como para generar predicciones agregadas útiles. Uno de los elementos más importantes es 

La teoría BLC postula desigualdades específicas en la propensión marginal a consumir según el origen de la riqueza, no arbitrarias, sino resultado de mecanismos adaptativos para facilitar el ahorro. Esta descomposición constituye un ejemplo de framing (Kahneman y Tversky, 1984). Mientras que la teoría tradicional asume invariancia al encuadre, el modelo BLC introduce dependencia del encuadre.

Para ilustrar el sistema de tres cuentas, consideremos un hogar que utiliza una regla de pensiones según la cual se deduce una fracción s de la renta en cada periodo, sin acceso a esos fondos hasta la jubilación. Los saldos serían:
	1.	Cuenta de renta corriente: I = (1 - s)y_t.
	2.	Cuenta de riqueza actual: $A = \sum_{i=1}^{t-1} \big[(1 - s)y_i - c_i\big]$.
	3.	Cuenta de riqueza futura: suma de ingresos futuros netos más la riqueza acumulada en el plan de pensiones.

Este esquema es una simplificación. En la práctica, los hogares subdividen aún más estas cuentas (por ejemplo, ahorro para educación). Además, existe ambigüedad en la clasificación de ciertos ingresos: los dividendos pueden tratarse como renta corriente, mientras que otras rentas del capital se asignan a activos. El patrimonio inmobiliario también puede clasificarse como riqueza futura o actual, según el hogar.

Aunque este sistema puede parecer extraño para un economista, es análogo al funcionamiento contable de muchas universidades privadas, que distinguen entre fondos corrientes y fondos patrimoniales (endowment), con reglas estrictas sobre su uso.

Si el individuo desea ahorrar más allá de lo permitido por el plan de pensiones, debe recurrir al ahorro discrecional, lo que requiere esfuerzo de autocontrol para generar ese ahorro adicional, no gastarlo antes de tiempo y evitar endeudarse contra ingresos futuros. El coste de ese esfuerzo depende inversamente de la tentación. En el modelo, la tentación de gastar un dólar adicional depende de su ubicación en las cuentas mentales: es mayor en la renta corriente, menor en los activos actuales y aún menor en la riqueza futura.



 La contraparte del autocontrol sería entonces la \textbf{tentación}, ya que algunas situaciones son menos propicias para el ahorro que otras















Apoyándonos en la investigación realizada sobre estos temas por psicólogos y otros científicos sociales (véanse, por ejemplo, Ainslie, 1975, y Mischel, 1981), somos capaces de formular predicciones específicas sobre cómo el comportamiento real de ahorro de los hogares difiere del modelo idealizado de ciclo vital.

Nuestro modelo produce tres proposiciones que están significativamente en conflicto con la hipótesis del ciclo vital en este ámbito general. En esta sección se analiza la sensibilidad del consumo a la renta. Las dos secciones siguientes se refieren a los casos especiales de bonificaciones y ganancias inesperadas (windfalls). 










\end{document}